\documentclass[11pt]{article}

%%% Packages
\usepackage{amsmath,amsthm,amssymb,amsfonts}
\usepackage{mathtools}
\usepackage{enumerate}
\usepackage{hyperref}
\usepackage{url}
\usepackage{booktabs}
\usepackage{array}
\usepackage{xcolor}
\usepackage{graphicx}
\usepackage{placeins}
\usepackage[margin=1in]{geometry}
\graphicspath{{../figures/}}

%%% Theorem environments
\newtheorem{theorem}{Theorem}[section]
\newtheorem{lemma}[theorem]{Lemma}
\newtheorem{proposition}[theorem]{Proposition}
\newtheorem{corollary}[theorem]{Corollary}
\newtheorem{conjecture}[theorem]{Conjecture}
\theoremstyle{definition}
\newtheorem{definition}[theorem]{Definition}
\newtheorem{example}[theorem]{Example}
\newtheorem{remark}[theorem]{Remark}
\newtheorem{observation}[theorem]{Observation}

%%% Notation shortcuts
\newcommand{\Z}{\mathbb{Z}}
\newcommand{\Q}{\mathbb{Q}}
\newcommand{\R}{\mathbb{R}}
\newcommand{\C}{\mathbb{C}}
\newcommand{\N}{\mathbb{N}}
\newcommand{\F}{\mathcal{F}}
\newcommand{\fpart}[1]{\left\{#1\right\}}
\DeclareMathOperator{\rank}{rank}
\DeclareMathOperator{\sgn}{sgn}

\title{The Geometric Signature of Primes in Farey Sequences}

\author{
  Saar Shai\thanks{Independent Researcher. Email: \texttt{saar.shai@gmail.com}.}
}

\date{April 2026}

\begin{document}

\maketitle

\begin{abstract}
We study how the uniformity of the Farey sequence changes
each time a new integer is included---a per-step perspective
that appears not to have been previously investigated.
Primes and composites play opposite roles: composites
improve uniformity while primes disrupt it.

\textbf{Proved results.}
We establish exact algebraic identities for the per-step
change~$\Delta W(p)$: a four-term decomposition into dilution,
cross, shift-squared, and new-fraction terms; an injection
principle showing each new prime fraction enters a distinct
gap; a dominance result for the fraction~$1/p$; and a
near-cancellation theorem showing the dilution and new-fraction
terms are equal to order~$O(1/p)$, so the net wobble change
is controlled by the cross and shift-squared terms alone.
All identities are formally verified in Lean~4
(441~results across 31~files; the Bridge Identity,
Injection Principle, four-term decomposition, and
Displacement-Shift identity are fully proved with
zero~\texttt{sorry} in their files).

\textbf{Computational observations.}
The sign of~$\Delta W(p)$ correlates ($r = 0.77$) with the
leading Riemann zero via the Mertens function and the
Chebyshev bias framework of Rubinstein--Sarnak.
A spectral function built from the per-step data detects
the first three nontrivial zeta zeros from Farey discrepancy
data alone.  A certified counterexample at $p = 92{,}173$
($M(p) = -2$, $\Delta W(p) > 0$) shows the na\"ive sign
conjecture fails; zero counterexamples are found among primes
with $M(p) \le -3$ up to $p = 100{,}000$.
\end{abstract}

\medskip
\noindent\textbf{Keywords:} Farey sequence, Mertens function, Chebyshev bias,
zeta zeros, rank deviation, Ramanujan sums, formal verification, Lean~4.

\noindent\textbf{MSC 2020:} 11B57, 11M06, 11N37, 11A25, 11Y35.

%% ================================================================
\section{Introduction}\label{sec:intro}
%% ================================================================

\subsection{Farey sequences and their discrepancy}

For a positive integer~$N$, the \emph{Farey sequence} $\F_N$ is the
ascending list of fractions $a/b$ with $0 \le a \le b \le N$ and
$\gcd(a,b) = 1$.  Its cardinality is
\begin{equation}
  |\F_N| = 1 + \sum_{k=1}^{N} \varphi(k).
\end{equation}

\begin{figure}[ht]
\centering
\includegraphics[width=0.55\textwidth]{fig_circle_farey.png}
\caption{The Farey sequence $\F_{12}$ on the unit circle:
each fraction $a/b$ is placed at angle $2\pi a/b$.  The
pattern is approximately uniform, with the largest gaps near
$0$ and $1$ (top of the circle) where the denominators
are smallest.}
\label{fig:farey-circle}
\end{figure}

\noindent Writing $n = |\F_N|$ and listing elements as
$f_0 < f_1 < \cdots < f_{n-1}$, the \emph{squared rank deviation}
(which we denote $W(N)$ and refer to as the \emph{wobble}; this is our terminology, not standard) is
\begin{equation}
  W(N) \;=\; \sum_{j=0}^{n-1} \left( f_j - \frac{j}{n} \right)^{\!2}.
\end{equation}

\begin{remark}[Natural normalization: $N \cdot W(N)$]
Since $W(N) \to 0$ as $N \to \infty$, the rate of decay is the
object of interest.  The Franel--Landau theorem states that RH
is equivalent to $W(N) = O(N^{-1+\varepsilon})$ for every
$\varepsilon > 0$, i.e., $N \cdot W(N) = O(N^\varepsilon)$.
The quantity $N \cdot W(N)$ removes the built-in decay and
isolates the true signal.  Computationally, $p \cdot W(p)$
grows slowly from $\approx 0.28$ at $p=2$ to $\approx 0.56$
at $p=113$ and $\approx 0.87$ at $p=2000$, so $W(N)$ decays
strictly slower than $1/N$.  The precise growth rate of
$N \cdot W(N)$---and whether it is $O(\log N)$,
$O((\log N)^\alpha)$ for some $0 < \alpha < 1$, or
something else---is an open problem even conditionally on RH.
\end{remark}

\begin{remark}[Interval vs.\ circle]
The Farey sequence lives on the \emph{closed interval} $[0,1]$,
with $f = 0$ and $f = 1$ as distinct elements.  Throughout this
paper we use circle diagrams (identifying $[0,1)$ with $\R/\Z$)
as a visualization aid, but the wobble $W(N)$
and all identities in Section~\ref{sec:identities} are stated
and proved on the interval, where the two boundary points
$f = 0$ and $f = 1$ contribute separately.  In particular,
the $+2$ in the bridge identity (Theorem~\ref{thm:bridge})
arises from these two distinct boundary contributions.
\end{remark}

\begin{figure}[ht]
\centering
\includegraphics[width=\textwidth]{fig_wobble_circle.png}
\caption{The wobble visualized on the circle.  Dark dots show
actual Farey fractions; red circles show ideal (equally spaced)
positions.  Red arcs connect each fraction to its ideal position---brighter arcs mean larger discrepancy.  \textbf{Left:}
$\F_5$ has high wobble (large gaps between actual and ideal).
\textbf{Center:} $\F_{12}$ has lower wobble (fractions closer
to uniform).  \textbf{Right:} adding prime $p = 13$ (red stars)
changes the ideal positions and the wobble.}
\label{fig:wobble-circle}
\end{figure}

\subsection{The Franel--Landau theorem}

Franel~\cite{Franel1924} and Landau~\cite{Landau1924} proved that
the Riemann Hypothesis is equivalent to the bound
$\sum_{j=0}^{n-1} |f_j - j/n| = O(N^{1/2+\varepsilon})$
for every $\varepsilon > 0$, where $n = |\F_N|$.
Equivalently (see Hardy--Wright~\cite{HardyWright2008}, Ch.~XVIII),
RH holds if and only if $W(N) = O(N^{-1+\varepsilon})$.
Thus the quantity~$W(N)$ defined above---which we call the
\emph{wobble} (our terminology for the squared rank
deviation)---encodes information about the non-trivial zeros
of the Riemann zeta function.
Recent work by Karvonen and
Zhigljavsky~\cite{KarvonenZhigljavsky2025} studies the maximum
mean discrepancy of Farey sequences, providing additional
quantitative connections between Farey uniformity and analytic
number theory.

\subsection{Motivating observation: primes fill voids uniformly}

Place the Farey fractions of $\F_N$ as points on the unit circle
at angles $2\pi f$.  As $N$ grows, new fractions appear and the
circle fills in.  The first author observed a striking asymmetry:
when a \emph{prime}~$p$ is added, every one of the $p-1$ new
fractions $k/p$ lands on the circle at perfectly equispaced
positions, filling voids uniformly across the entire circle.
This happens because every integer below a prime is coprime to
it ($\varphi(p) = p-1$), so no position is skipped.  Composites,
by contrast, skip positions whose numerators share a factor with
the denominator, leaving voids in some regions while clustering
in others.

Moreover, by the standard mediant property of Farey sequences,
each new fraction $k/p$ is the \emph{mediant} of its two
neighbors in $\F_{p-1}$---the geometrically optimal insertion
point.  This means primes do not merely add points; they add
them in the most balanced way possible.

These geometric observations led to a natural question:
does the uniformity of the Farey sequence on the circle
improve or worsen at each step, and what controls the answer?

\begin{figure}[ht]
\centering
\includegraphics[width=\textwidth]{fig_void_filling.png}
\caption{How primes break voids.  \textbf{Left:} $\F_6$ with
voids highlighted by size (darker = larger).
\textbf{Center:} prime $p = 7$ inserts one fraction into each of
$p - 1 = 6$ distinct gaps (green arcs), filling some voids and
skipping others.
\textbf{Right:} composite $N = 8$ leaves voids unfilled (red
arcs) due to missing coprime numerators.}
\label{fig:voids}
\end{figure}

\subsection{Per-step decomposition: a new perspective}

Despite decades of work on the Franel--Landau framework, the
\emph{per-step} behavior---how $W$ changes when a single integer
is appended---appears never to have been studied.
Steinerberger~\cite{Steinerberger2019} constructs \emph{greedy}
sequences on $[0,1]$ that minimize a logarithmic potential energy
at each step, achieving discrepancy $D_N \lesssim (\log N)/\sqrt{N}$
and conjecturally the Schmidt-optimal $O((\log N)/N)$.
The Farey construction is fundamentally different: it is
defined by arithmetic (coprimality), not by discrepancy
minimization.  As we show, prime levels~$\F_p$ always insert
new fractions at their mediant positions---the local geometric
optimum within each gap---yet this can \emph{increase} global
discrepancy.  This contrast illustrates why the per-step view
is nontrivial: geometric optimality at the local level does not
imply discrepancy minimization at the global level.  We define
\begin{equation}
  \Delta W(N) \;=\; W(N{-}1) - W(N),
\end{equation}
so $\Delta W(N) > 0$ means adding $N$ \emph{improves} uniformity,
while $\Delta W(N) < 0$ means it \emph{worsens} it.

The per-step view reveals a sharp asymmetry invisible in the
cumulative $W(N)$: composites almost always improve uniformity
($96\%$ of composites $N \le 200$ have $\Delta W > 0$), while
primes are the primary source of ``damage'' ($99\%$ of
negative-$\Delta W$ events come from primes, in computations
through $N = 100{,}000$).  Moreover, for
primes, the sign of $\Delta W(p)$ tracks the Mertens function
$M(p) = \sum_{k \le p} \mu(k)$---the Mertens function acts as
a ``dial'' controlling the prime contributions.  This creates
a direct bridge between two worlds that had no known connection
at this level of granularity: the geometry of Farey discrepancy
and the arithmetic of the M\"obius function.

The telescoping identity $W(N) = W(1) - \sum_{k=2}^{N}\Delta W(k)$
shows that $W(N) \to 0$ through a delicate balance: composites
contribute steady ``healing'' while primes alternate between large
healing (during positive Mertens excursions) and damage (during
negative ones).  This decomposition of Farey convergence into
prime and composite contributions was previously unknown.

\begin{figure}[ht]
\centering
\includegraphics[width=\textwidth]{farey_spectroscope.png}
\caption{The Farey spectroscope (Section~\ref{sec:spectroscope}):
the spectral function
$F(\gamma) = |\sum R(p)\,p^{-1/2-i\gamma}|^2$ (where $R(p)$ is the
correlation ratio defined in Section~\ref{sec:spectroscope}),
built entirely from per-step Farey discrepancy data, exhibits sharp peaks at
the locations of nontrivial zeta zeros (dashed lines).  The
first three zeros $\gamma_1, \gamma_2, \gamma_3$ are recovered
to within~$0.8\%$.}
\label{fig:spectroscope-preview}
\end{figure}

%% ================================================================
\section{Definitions and Notation}\label{sec:defs}
%% ================================================================

\begin{definition}[Rank discrepancy]
For $f_j \in \F_N$ with $n = |\F_N|$, the \emph{rank discrepancy} is
$D(f_j) = j - n \cdot f_j$.
\end{definition}

\begin{remark}[Relation to local discrepancy]
Our $D(f)$ equals $-\hat{r}_N(f)$ in the notation of
Niederreiter~\cite{HardyWright2008} and
Tomas~Garcia~\cite{TomasGarcia2025}, who define the local
discrepancy as $\hat{r}_N(f) = f\,|\F_N| - I_N(f)$
(expected rank minus actual).  The sign convention difference
is standard in the literature; we adopt the ``actual minus
expected'' convention throughout.
\end{remark}

\begin{definition}[Shift]
For a prime~$p$ and $f \in \F_{p-1}$:
$\delta(f) = f - \fpart{pf}$,
where $\fpart{x} = x - \lfloor x \rfloor$ is the fractional part.
\end{definition}

\begin{definition}[Per-denominator shift sum]
For denominator~$b$ with $2 \le b \le N$, define
$\delta_b^2 = \sum_{\substack{a=1 \\ \gcd(a,b)=1}}^{b} \delta(a/b)^2$,
the total squared shift contributed by fractions with denominator~$b$.
\end{definition}

\begin{definition}[Mertens function]
$M(N) = \sum_{k=1}^{N} \mu(k)$,
where $\mu$ is the M\"obius function.
\end{definition}


%% ================================================================
\section{New Identities}\label{sec:identities}
%% ================================================================

The identities in this section follow from classical
Ramanujan-sum decompositions~\cite{Ramanujan1918, Kanemitsu1982}
combined with the Farey involution and the M\"obius function.
The evaluation $\sum_{f \in \F_N} e^{2\pi i f} = M(N)+1$ is
implicit in the Farey-series chapter of Edwards~\cite{Edwards1974}
(see also Hardy--Wright~\cite{HardyWright2008}, Ch.~XVI\@).

The per-frequency generalization (Theorems~\ref{thm:gen-bridge}
and~\ref{thm:universal}) and the per-step application to
$\Delta W(N)$ (Theorems~\ref{thm:bridge}--\ref{thm:cross-term})
appear, to our knowledge, not to have been stated previously
in this per-step form.  (The underlying exponential sums are
Ramanujan sums, whose evaluation is classical~\cite{Ramanujan1918};
our contribution is the per-step framework connecting them
to~$\Delta W$.)  All identities and supporting lemmas are
formally verified in Lean~4 (422~results across fifteen files,
zero~\texttt{sorry}); all prime-parameter identities
are also verified
computationally for all primes $p \le 100{,}000$, and the universal
formula (Theorem~\ref{thm:universal}) is verified on the grid
$m \le 40$, $N \le 30$.

\subsection{The Bridge Identity}

\begin{theorem}[Bridge Identity]\label{thm:bridge}
For every prime $p \ge 2$,
\begin{equation}
  \sum_{f \in \F_{p-1}} e^{2\pi i p f}
  \;=\; M(p) + 2.
\end{equation}
\end{theorem}

\begin{proof}
Decompose by denominator.  The boundary terms $f = 0$ and $f = 1$
contribute $1 + 1 = 2$.  For each denominator $2 \le b \le p-1$,
the inner sum over numerators coprime to $b$ is the Ramanujan sum
$c_b(p)$.  Since $\gcd(p, b) = 1$ for $b < p$, we have
$c_b(p) = \mu(b)$.  Summing: $2 + \sum_{b=2}^{p-1} \mu(b) =
2 + M(p-1) - 1 = M(p) + 2$, using $M(p-1) = M(p) + 1$.
\end{proof}

\begin{remark}
Taking real parts gives the cosine form
$\sum \cos(2\pi pf) = M(p) + 2$;
the imaginary part vanishes by the symmetry $f \leftrightarrow 1 - f$.
\end{remark}

The bridge identity is a special case of a more general result:

\begin{theorem}[Generalized Bridge Identity]\label{thm:gen-bridge}
For any positive integer $N$ and any prime $m > N$,
\begin{equation}
  \sum_{f \in \F_N} e^{2\pi i m f} \;=\; M(N) + 1.
\end{equation}
In particular, all primes $m > N$ yield the \emph{same}
exponential sum over $\F_N$, depending only on $M(N)$.
\end{theorem}

\begin{proof}
Same decomposition as Theorem~\ref{thm:bridge}.  Every
denominator $b \le N < m$ satisfies $\gcd(b, m) = 1$,
so $c_b(m) = \mu(b)$.  The boundary gives~$2$; the
interior gives $M(N) - 1$.  Total: $M(N) + 1$.
\end{proof}

\begin{remark}
Theorem~\ref{thm:bridge} is the case $N = p - 1$:
$M(p-1) + 1 = (M(p) + 1) + 1 = M(p) + 2$.
For tensor-product Farey sets in $d$~dimensions, the sum
at a prime-vector frequency is $(M(N)+1)^d$; when $M(N) = -1$
(e.g., $N = 10$), this vanishes in every dimension.
\end{remark}

Both Theorems~\ref{thm:bridge} and~\ref{thm:gen-bridge} are
special cases of a universal formula:

\begin{theorem}[Universal Farey Exponential Sum]\label{thm:universal}
For any positive integers $m$ and $N$,
\begin{equation}
  \sum_{f \in \F_N} e^{2\pi i m f}
  \;=\; M(N) + 1
  + \sum_{\substack{d \mid m \\ d \ge 2}} d \cdot M\!\left(\left\lfloor N/d \right\rfloor\right).
\end{equation}
\end{theorem}

\begin{proof}
Group the sum by denominator:
\begin{equation*}
\sum_{f \in \F_N} e^{2\pi i m f}
= \sum_{b=1}^{N} \sum_{\substack{a=0 \\ \gcd(a,b)=1}}^{b}
e^{2\pi i m a/b}.
\end{equation*}
The boundary fractions $a/b \in \{0/1,\, 1/1\}$ contribute
$e^0 + e^{2\pi im} = 1 + 1 = 2$.
For $b \ge 2$, the inner sum over residues $a$ with
$0 \le a < b$ and $\gcd(a,b)=1$ is the Ramanujan sum
$c_b(m) = \sum_{\substack{a=1\\ \gcd(a,b)=1}}^{b} e^{2\pi i m a/b}$
\cite{Ramanujan1918}.  Thus:
\begin{equation*}
\sum_{f \in \F_N} e^{2\pi i m f} = 2 + \sum_{b=2}^{N} c_b(m).
\end{equation*}

Apply the Möbius evaluation of the Ramanujan sum
(\cite{HardyWright2008}, Thm.~272):
\begin{equation*}
c_b(m) = \sum_{d \mid \gcd(b,m)} d\,\mu(b/d).
\end{equation*}
Writing $b = dk$ and exchanging summation order,
\begin{equation*}
\sum_{b=2}^{N} c_b(m)
= \sum_{b=2}^{N}\ \sum_{d \mid \gcd(b,m)} d\,\mu(b/d)
= \sum_{\substack{d \mid m}} d \sum_{\substack{k=1 \\ dk \le N}}^{\lfloor N/d \rfloor} \mu(k).
\end{equation*}
The interchange is valid since every pair $(b,d)$ with $d \mid \gcd(b,m)$,
$2 \le b \le N$, corresponds uniquely to $(d,k)$ with $d \mid m$, $k = b/d \ge 1$,
$dk \le N$.

Separate the $d=1$ contribution:
\begin{equation*}
\sum_{\substack{d \mid m \\ d \ge 1}} d \sum_{k \le \lfloor N/d \rfloor} \mu(k)
= \sum_{k=1}^{N} \mu(k) + \sum_{\substack{d \mid m \\ d \ge 2}} d \cdot M\!\left(\lfloor N/d \rfloor\right).
\end{equation*}
The first sum is $M(N)$.  Adding the boundary term~$2$ and
subtracting the $k=1$ term (since the $b$-sum starts at $b=2$,
i.e.\ $k \ge 1$ but $d=1, k=1$ gives $b=1$ which is excluded):
more precisely, for $d=1$ the sum runs over $k=2,\ldots,N$, giving
$\sum_{k=2}^{N}\mu(k) = M(N)-1$.  Hence:
\begin{equation*}
2 + \sum_{b=2}^{N} c_b(m)
= 2 + (M(N) - 1) + \sum_{\substack{d \mid m \\ d \ge 2}} d\cdot M\!\left(\lfloor N/d \rfloor\right)
= M(N) + 1 + \sum_{\substack{d \mid m \\ d \ge 2}} d\cdot M\!\left(\lfloor N/d \rfloor\right). \qedhere
\end{equation*}
\end{proof}

\begin{remark}
This is a \emph{complete} evaluation of the Farey exponential sum
at every integer frequency in terms of the Mertens function at
divisor-scaled arguments.  For prime $m > N$, the only divisor
$d \ge 2$ is $m$ itself, and $\lfloor N/m \rfloor = 0$ gives
$M(0) = 0$, recovering Theorem~\ref{thm:gen-bridge}.
Verified computationally for all $m \le 40$ and $N \le 30$.
\end{remark}

\begin{figure}[ht]
\centering
\includegraphics[width=\textwidth]{fig_universal_formula.png}
\caption{The universal formula for $\F_{20}$.
\textbf{Top:} The Farey exponential sum at each frequency $m$.
Green bars (primes $m > 20$) all equal $M(20)+1 = -2$---they
are \emph{identical}, independent of which prime.  Blue bars
(primes $m \le 20$) and orange bars (composites) vary.
\textbf{Bottom:} prime indicator.  The dashed line at
$m = N = 20$ separates the two regimes.}
\label{fig:universal}
\end{figure}

\begin{figure}[ht]
\centering
\includegraphics[width=\textwidth]{fig_bridge_vectors.png}
\caption{The Bridge Identity visualized.  Each Farey fraction
$f \in \F_{p-1}$ becomes a point at angle $2\pi pf$ on the
unit circle.  Their vector sum (red arrow) equals $M(p) + 2$.
\textbf{Left:} $p = 11$, $M = -2$, sum $= 0$ (vectors cancel).
\textbf{Right:} $p = 97$, $M = 1$, sum $= 3$ (net rightward bias).}
\label{fig:bridge-vectors}
\end{figure}


\subsection{The Displacement--Cosine Identity}

\begin{theorem}[Displacement--Cosine Identity]\label{thm:disp-cos}
For every odd prime $p$,
\begin{equation}
  \sum_{f \in \F_{p-1}} D(f) \cdot \cos(2\pi p f)
  \;=\; -1 - \frac{M(p)}{2}.
\end{equation}
\end{theorem}

\begin{proof}
Partition $\F_{p-1}$ into four parts.

\smallskip\noindent\textit{Boundary $f=0$.}\quad
$D(0) = 0$, so the contribution is~$0$.

\smallskip\noindent\textit{Boundary $f=1$.}\quad
$D(1) = (n-1) - n \cdot 1 = -1$ and $\cos(2\pi p) = 1$,
contributing~$-1$.

\smallskip\noindent\textit{Midpoint $f = 1/2$.}\quad
Since $n = |\F_{p-1}|$ is always odd (the symmetry
$f \leftrightarrow 1-f$ pairs all non-boundary, non-midpoint
elements, so $n = 2k + 3$ for some $k \ge 0$),
$\operatorname{rank}(1/2) = (n-1)/2$ and
$D(1/2) = (n-1)/2 - n/2 = -1/2$.
For odd prime $p$, $\cos(\pi p) = -1$, contributing
$(-1/2)(-1) = 1/2$.

\smallskip\noindent\textit{Interior pairs $\{f,\, 1-f\}$ with $f \in (0,1/2)$.}\quad
By the antisymmetry proved in the cross-term section
(Theorem~\ref{thm:cross-term}), $D(f) + D(1-f) = -1$.
Since $\cos(2\pi p(1-f)) = \cos(2\pi p - 2\pi pf) = \cos(2\pi pf)$
(integer period), each pair contributes
$[D(f) + D(1-f)] \cdot \cos(2\pi pf) = -\cos(2\pi pf)$.

\smallskip\noindent Summing all parts:
\begin{equation*}
  \sum_{f \in \F_{p-1}} D(f)\cos(2\pi pf)
  = 0 + \tfrac{1}{2} + (-1) - S
  = -\tfrac{1}{2} - S,
\end{equation*}
where $S = \sum_{f \in (0,1/2) \cap \F_{p-1}} \cos(2\pi pf)$.

\smallskip\noindent\textit{Evaluating $S$ via the Bridge Identity.}\quad
Taking the real part of the Bridge Identity
(Theorem~\ref{thm:bridge}):
$\sum_{f \in \F_{p-1}} \cos(2\pi pf) = M(p) + 2$.
Decomposing this sum as $1$ (from $f=0$) $+$ $(-1)$ (from
$f = 1/2$, since $\cos(\pi p) = -1$) $+$ $1$ (from $f=1$)
$+$ $2S$ (pairing interior fractions), we get
$1 - 1 + 1 + 2S = M(p) + 2$, hence
$S = (M(p)+1)/2$.

\smallskip\noindent Substituting:
$-1/2 - S = -1/2 - (M(p)+1)/2 = -1 - M(p)/2$.
\end{proof}

\begin{remark}
This is the keystone identity: the Fourier coefficient of the
rank discrepancy at frequency $p$ is a \emph{linear} function
of the Mertens function.  It directly connects how far each
Farey fraction sits from its ideal position (geometry) to the
cancellations in the M\"obius function (arithmetic).
\end{remark}



\subsection{Compact Cross-Term Formula}

\begin{theorem}\label{thm:cross-term}
For every prime $p \ge 2$,
\begin{equation}
  \sum_{f \in \F_{p-1}} D(f) \cdot \delta(f)^2
  = -\frac{1}{2}\sum_{f \in \F_{p-1}} \delta(f)^2 - \frac{1}{2}.
\end{equation}
\end{theorem}

\begin{proof}
Let $n = |\F_{p-1}|$ and write $\Sigma = \sum_{f \in \F_{p-1}} D(f)\cdot\delta(f)^2$.
Using the convention $D(f_j) = j - n \cdot f_j$ with
zero-based index $j = 0, 1, \ldots, n-1$ (so $f_0 = 0$ and
$f_{n-1} = 1$), we verify the two key facts:

\medskip
\noindent\textit{Boundary values.}
$D(0) = 0 - n\cdot 0 = 0$ and $\delta(0) = 0 - \{0\} = 0$, so $f_0=0$
contributes $0$.
$D(1) = (n-1) - n\cdot 1 = -1$ and $\delta(1) = 1 - \{p\} = 1 - 0 = 1$,
so $f_{n-1}=1$ contributes $D(1)\delta(1)^2 = -1$.

\medskip
\noindent\textit{Antisymmetry of $D$ and $\delta$.}
For interior $f = a/b \in (0,1)$ with $b \le p-1$, the fraction
$1-f = (b-a)/b$ is also interior and in $\F_{p-1}$.
If $f$ has rank $j$ (zero-based) then by Farey symmetry
$1-f$ has rank $n-1-j$, giving
\begin{equation*}
D(f) + D(1-f) = j - nf + (n-1-j) - n(1-f) = -1.
\end{equation*}
For $\delta$: since $b < p$ implies $\gcd(b,p)=1$, the value
$pf = pa/b$ is not an integer, so $\{p(1-f)\} = 1 - \{pf\}$.
Hence $\delta(1-f) = (1-f) - (1-\{pf\}) = \{pf\} - f = -\delta(f)$,
and in particular $\delta(1-f)^2 = \delta(f)^2$.

\medskip
\noindent\textit{Pairing the interior.}
Note that $\delta(1/2) = 1/2 - \{p/2\} = 0$ for odd~$p$
(since $p/2 = k + 1/2$ has fractional part $1/2$),
so the midpoint $1/2 \in \F_{p-1}$ contributes~$0$.
For each interior pair $\{f, 1-f\}$ with $f \in (0, 1/2)$:
\begin{equation*}
D(f)\delta(f)^2 + D(1-f)\delta(1-f)^2
= \bigl[D(f) + D(1-f)\bigr]\delta(f)^2 = -\delta(f)^2.
\end{equation*}
Summing over the $\frac{n-3}{2}$ such pairs (excluding the midpoint)
plus the boundary contribution $D(1)\delta(1)^2 = -1$:
\begin{equation*}
\Sigma = -\sum_{\substack{f \in (0,1)}} \delta(f)^2\cdot\tfrac{1}{2} + (-1).
\end{equation*}
(The factor $\frac{1}{2}$ accounts for each pair appearing once.)
Since $\delta(0)^2 = 0$, we may include $f=0$ freely; and
$\delta(1)^2 = 1$, so
$-\tfrac{1}{2}\sum_{f\ne 0,1}\delta^2 - 1
= -\tfrac{1}{2}\!\left(\sum_f \delta^2 - 0 - 1\right) - 1
= -\tfrac{1}{2}\sum_f \delta^2 + \tfrac{1}{2} - 1
= -\tfrac{1}{2}\sum_f \delta^2 - \tfrac{1}{2}$.
\end{proof}


\subsection{Geometric Identity: $B+C$ as $L^2$ Norm Change}

Recall $\delta(f) = \lfloor pf \rfloor - (p-1)f$ is the per-fraction shift
from Definition~\ref{def:shift} and $D(f)$ is the displacement.
The cross term $B = (2/{n'}^2)\sum_{f \in \F_{p-1}} D(f)\cdot\delta(f)$ and the
shift-squared term $C = (1/{n'}^2)\sum_{f \in \F_{p-1}} \delta(f)^2$
(see the full four-term decomposition in Section~\ref{sec:phenomenon}) admit a
geometric interpretation.

\begin{lemma}[Geometric Identity]\label{thm:geometric-bc}
For every prime $p \ge 2$,
\begin{equation}
  B + C \;=\; \frac{1}{{n'}^2}\!\left(
    \sum_{f \in \F_{p-1},\, f \ne 1}
      \bigl(D_{\F_{p-1}}(f) + \delta(f)\bigr)^2
    \;-\; \sum_{f \in \F_{p-1},\, f \ne 1} D_{\F_{p-1}}(f)^2
  \right),
\end{equation}
where $n' = |\F_p|$.  In particular,
$B + C > 0$ if and only if adding $\delta$ to $D$ increases
the $L^2$ norm of the displacement over old fractions
(excluding $f = 1$).
\end{lemma}

\begin{proof}
Expanding $(D + \delta)^2 - D^2 = 2D\delta + \delta^2$ and
summing gives $2\sum D\delta + \sum \delta^2$, which equals
$n'^2(B + C)$ by the definitions of $B$ and $C$.
\end{proof}

\begin{remark}
The physical interpretation is immediate: inserting a prime
perturbs each old fraction's displacement by $\delta(f)$
(Proposition~\ref{prop:disp-shift}).  The identity says that
$B + C > 0$ is
equivalent to the statement that this perturbation \emph{increases}
the total spectral disorder (squared displacement).
Computationally, $B + C > 0$ holds for all primes
$p \ge 11$ with $M(p) \le -3$ (tested up to $p = 100{,}000$),
but can fail for primes with large positive Mertens values
(e.g.\ $B + C < 0$ at $p = 1399$ with $M(p) = +8$).
The Sign Pattern does not depend on $B + C > 0$:
it uses the stronger bypass $C + \mathcal{D} > A$
(see Section~\ref{sec:sign-theorem}).
This is formally verified in Lean~4 (\texttt{bPlusC\_eq\_new\_minus\_old\_sq}).
\end{remark}

\begin{lemma}[Permutation Square-Sum Identity]\label{thm:perm-sq}
For every prime $p$ and every $N < p$,
\begin{equation}
  \sum_{\substack{f = a/b \in \F_N \\ b > 1}}
    f \cdot \delta(f)
  \;=\; \frac{1}{2}
  \sum_{\substack{f = a/b \in \F_N \\ b > 1}}
    \delta(f)^2.
\end{equation}
Equivalently, using the notation of the four-term decomposition,
$\sum f\,\delta(f) = \tfrac{1}{2}\,{n'}^2 C$.
\end{lemma}

\begin{proof}
For $f = a/b$ with $b \ge 2$ and $\gcd(a,b)=1$, write
$\delta(f) = f - \{pf\}$ where $\{pf\} = (pa \bmod b)/b$
denotes the fractional part.  Then:
\begin{equation*}
f\,\delta(f) - \tfrac{1}{2}\delta(f)^2
= f(f - \{pf\}) - \tfrac{1}{2}(f-\{pf\})^2
= \tfrac{1}{2}(f^2 - \{pf\}^2).
\end{equation*}
Thus it suffices to show
\begin{equation*}
\sum_{\substack{f = a/b \in \F_N \\ b \ge 2}} (f^2 - \{pf\}^2) = 0.
\end{equation*}
Fix denominator $b \ge 2$ with $b \le N$.
Since $p < b$ is impossible for $b \le N < p$, we have
$\gcd(p, b) = 1$.  Therefore the map
$\sigma_b : a \mapsto pa \bmod b$
is a bijection on $\{a : 1 \le a \le b-1,\, \gcd(a,b)=1\}$
(multiplication by $p$ is invertible mod $b$).
Hence:
\begin{equation*}
\sum_{\substack{a=1 \\ \gcd(a,b)=1}}^{b-1} \left(\frac{a}{b}\right)^2
= \sum_{\substack{a=1 \\ \gcd(a,b)=1}}^{b-1} \left(\frac{pa \bmod b}{b}\right)^2.
\end{equation*}
Subtracting, each per-denominator contribution to
$\sum (f^2 - \{pf\}^2)$ is zero.  Summing over all
$b = 2, \ldots, N$ gives $0$, whence
\begin{equation*}
\sum_{\substack{f \in \F_N \\ b \ge 2}} f\,\delta(f)
= \tfrac{1}{2}\sum_{\substack{f \in \F_N \\ b \ge 2}} \delta(f)^2. \qedhere
\end{equation*}
\end{proof}

\begin{remark}
Combined with the Farey rank decomposition
$D(f) = -f - R(f)$, where $R(f) = \sum_{d \le N} \mu(d)\sum_{m=1}^{\lfloor N/d\rfloor} \{fm\}$
arises from M\"obius inversion of the counting function,
identity~the permutation identity yields
\[
  B + C \;=\; \frac{-2}{{n'}^2}\sum_{f \in \F_N,\, b > 1} R(f)\,\delta(f).
\]
Thus $B + C > 0$ is equivalent to a \emph{negative} correlation
between the M\"obius-weighted fractional-part function~$R$
and the multiplicative displacement~$\delta$.
For primes with $M(p) \le -3$, this correlation is
empirically always negative; the sign reversal at primes
with large positive~$M(p)$ reflects the Mertens--discrepancy
connection (Section~\ref{sec:phenomenon}).
\end{remark}

\subsection{Displacement-Shift Identity}

The following proposition explains \emph{why} the shift
function~$\delta$ controls the per-step wobble change.

\begin{proposition}[Displacement-Shift]\label{prop:disp-shift}
In the transition from $\F_{p-1}$ to $\F_p$ (where $p$ is prime),
the displacement of each old fraction $f \in \F_{p-1}$, $f \ne 1$,
changes by exactly the shift:
\begin{equation}
  D_{\F_p}(f) \;=\; D_{\F_{p-1}}(f) + \delta(f).
\end{equation}
At the boundary $f = 1$: $D_{\F_p}(1) = D_{\F_{p-1}}(1)$.
\end{proposition}

\begin{proof}
The rank of $f = a/b$ in $\F_p$ equals its rank in $\F_{p-1}$ plus the
number of fractions $k/p$ inserted below~$f$, which is
$\lfloor pf \rfloor$.  Thus
$D_{\F_p}(f) = (\rank_{\F_{p-1}}(f) + \lfloor pf \rfloor) -
(n + p - 1) \cdot f = D_{\F_{p-1}}(f) + \lfloor pf \rfloor -
(p-1)f = D_{\F_{p-1}}(f) + (f - \{pf\}) = D_{\F_{p-1}}(f) +
\delta(f)$.  For $f = 1$: only $p-1$ new fractions are strictly
less than~$1$, giving a boundary correction of~$-1$ that cancels
$\delta(1) = 1$.
\end{proof}

\begin{remark}
This identity makes precise the mechanism by which the Mertens
function controls $\Delta W$: the cross term
$\sum D_{\F_{p-1}}(f) \cdot \delta(f)$---which our other identities
partially evaluate---is exactly the linear interaction between the
old pattern and the shift caused by inserting the new prime.
\end{remark}

\begin{proposition}[Zero-sum at equispaced points]

For any prime~$p \ge 3$,
$\sum_{k=1}^{p-1} D_{\F_{p-1}}(k/p) = 0$.
\end{proposition}

\begin{proof}
Write $S = \sum_{k=1}^{p-1} D_{\F_{p-1}}(k/p)
= \sum_{k=1}^{p-1} \rank(k/p, \F_{p-1}) - n \cdot (p-1)/2$.
The sum $\sum_{k=1}^{p-1} \rank(k/p, \F_{p-1})$
counts pairs $(k, f)$ with $f \in \F_{p-1}$, $f \le k/p$.
Exchanging summation order, each $f = a/b$ contributes
$p - \lceil pa/b \rceil$ terms
(with the convention $\lceil 0 \rceil = 0$ at $f=0$,
contributing $p-1$ terms since $k$ ranges over $1,\ldots,p{-}1$).  By Farey symmetry
($f \leftrightarrow 1{-}f$), the sum
$\sum_{f \in \F_{p-1}} \lceil pf \rceil = n(p+1)/2 - 1$.
Hence $\sum \rank = (np - 1) - (n(p+1)/2 - 1) = n(p-1)/2$,
giving $S = 0$.
Verified independently by exact rational arithmetic for
all primes $p \le 500$.
\end{proof}


\subsection{Simultaneous Mertens extraction (``Mertens tomography'')}

The inversion of the universal formula provides a multi-scale
view of the Mertens function from a single Farey
sequence. We refer to this as \emph{Mertens tomography} (our terminology).

\begin{remark}[Mertens at multiple scales]
The universal formula at prime frequency~$p \le N$ gives
$M(\lfloor N/p \rfloor) = [S(p,N) - S(1,N)]/p$,
where $S(m,N) = \sum_{f \in \F_N} e^{2\pi imf}$.
Evaluating at multiple primes from a single $\F_N$ simultaneously
recovers $M$ at all scales $\lfloor N/p \rfloor$; composite
frequencies provide redundant consistency checks via
the universal formula.
\end{remark}


%% ================================================================
\section{The Wobble--Mertens Phenomenon}\label{sec:phenomenon}
%% ================================================================

\subsection{Empirical observation}

Computation through $p = 100{,}000$ reveals a striking pattern:
among primes with $\Delta W(p) > 0$ (wobble decreasing, i.e.,
uniformity improving), essentially all have $M(p) \ge 0$.  Among
primes with $M(p) < 0$, essentially all have $\Delta W(p) \le 0$
(uniformity worsening or unchanged).

\begin{center}
\begin{tabular}{lrrr}
\toprule
Range & Primes $p \ge 11$ & $\Delta W > 0$ & $M < 0$ with $\Delta W > 0$ \\
\midrule
$p \le 1{,}000$ & $164$ & $0$ (0.0\%) & 0 \\
$p \le 5{,}000$ & $665$ & $43$ (6.5\%) & 0 \\
$p \le 10{,}000$ & $1{,}225$ & $110$ (9.0\%) & 0 \\
$p \le 20{,}000$ & $2{,}258$ & $342$ (15.1\%) & 0 \\
$p \le 50{,}000$ & $5{,}129$ & $1{,}109$ (21.6\%) & 0 \\
$p \le 100{,}000$ & $9{,}588$ & $2{,}295$ (23.9\%) & 1 \\
\bottomrule
\end{tabular}
\end{center}

\noindent The positive-$\Delta W$ rate (proportion of primes with
$\Delta W > 0$, which we call ``violations'' of the general
downward trend) is zero below $p \approx 1{,}400$ because
$\Delta W$ is large and negative for small primes regardless
of the sign of~$M$.  Positive $\Delta W$ first appears near
$p = 1{,}399$ as $M(p)$ enters
its first sustained positive excursion, and the rate oscillates
between 15--27\% thereafter, tracking the Mertens function.

Binning primes by their normalized Mertens value $M(p)/\sqrt{p}$
reveals a sharp sigmoid relationship:

\begin{center}
\begin{tabular}{lr}
\toprule
$M/\sqrt{p}$ range & Positive-$\Delta W$ rate \\
\midrule

\end{tabular}
\end{center}
% [Detailed empirical tables available in supplementary]

\subsection{The threshold question}

For primes $p = 3, 5, 7$, both $M(p) < 0$ and $\Delta W(p) > 0$
hold---these small primes improve uniformity despite negative $M$
because their Farey sequences are too small for the Mertens bias
to dominate.  Starting from $p = 11$, the pattern reverses:
computational evidence for $p \le 50{,}000$ shows zero exceptions
among $5{,}129$ primes with $p \ge 11$.  A natural question is
whether $M(p) < 0$ alone suffices to guarantee $\Delta W(p) \le 0$.
It does not.

\begin{observation}[Certified counterexample at $p = 92{,}173$]
At the prime $p = 92{,}173$:
\[
  M(92{,}173) = -2, \qquad
  \Delta W(92{,}173) = +3.56 \times 10^{-11}.
\]
The Mertens function is strictly negative, yet the wobble decreases.
This is confirmed by four independent computations:
\begin{enumerate}[(i)]
  \item Naive double-precision floating point.
  \item Long double (80-bit extended precision).
  \item Kahan compensated summation (double precision with error tracking).
  \item \textbf{256-bit MPFR arithmetic} (high-precision numerical verification):
    $W(92{,}172) - W(92{,}173) = +3.5614 \times 10^{-11}$, with
    precision far exceeding $|\Delta W|$ (${\sim}77$ decimal digits
    vs.\ $|\Delta W| \approx 10^{-11}$).
\end{enumerate}
The sign is confirmed beyond reasonable numerical doubt.
This is the \emph{only} counterexample among the $9{,}588$ primes
$p \ge 11$ up to $p = 100{,}000$ (the small primes $p = 3, 5, 7$
are excluded; see Section~\ref{sec:phenomenon}).
\end{observation}


\subsection{The $M \le -3$ class: zero counterexamples}

Restricting to the stronger condition $M(p) \le -3$ eliminates
the counterexample entirely.

\begin{observation}[Computational: zero counterexamples for $M \le -3$]
Among all $4{,}617$ primes $p \le 100{,}000$ with $M(p) \le -3$:
\[
  \Delta W(p) \le 0 \quad\text{in every case.}
\]
Zero counterexamples.  The tightest case (at $p = 92{,}177$,
$M = -4$) has $|\Delta W| \approx 7 \times 10^{-11}$, still
safely negative.
\end{observation}

The data supports a density-one version of the sign pattern:
$\Delta W(p) \le 0$ holds for all $4{,}617$ primes $p \le 100{,}000$
with $M(p) \le -3$, but a counterexample to the universal statement
exists at $p = 243{,}799$ (found in subsequent computation).
We therefore conjecture that the sign pattern $\sgn(\Delta W(p)) = \sgn(-M(p))$
holds for a \emph{density-one} subset of primes with $M(p) \le -3$,
analogous to the Chebyshev bias in prime races~\cite{RubinsteinSarnak1994}.
The counterexample at $p = 92{,}173$ (with $M = -2$) shows the
threshold $M_0 = -3$ is necessary.

\begin{figure}[ht]
\centering
\includegraphics[width=\textwidth]{fig_delta_w_signs.png}
\caption{Per-prime wobble change $\Delta W(p)$ for all primes
up to $p = 20{,}000$.  Red bars: $\Delta W > 0$ (uniformity
improves).  Blue bars: $\Delta W < 0$ (uniformity worsens).
The burst--quiet pattern is visible: positive-$\Delta W$
episodes cluster in bands corresponding to positive Mertens
excursions.}
\label{fig:delta-w}
\end{figure}

\subsection{The four-term decomposition of $\Delta W$}

The displacement-shift identity (Proposition~\ref{prop:disp-shift})
yields an exact algebraic decomposition of $\Delta W(p)$ into four
terms.

\begin{equation}
  \Delta W(p) = A - B - C - \mathcal{D}
\end{equation}
where $n = |\F_{p-1}|$, $n' = |\F_p| = n + (p-1)$, and
$B$, $C$ are summed over old fractions \emph{excluding}
$f = 1$.  (An earlier version included boundary correction
terms $+1 - 1/{n'}^2$; these arise only when $f = 1$ is
included in $B$ and~$C$.  With $f = 1$ excluded from both
sums, no boundary correction is needed.  Both conventions
give the same value of~$\Delta W$; the version without
boundary terms is cleaner.)
The decomposition is formally verified
in Lean~4 (\texttt{four\_term\_decomposition}).  The
terms are
\begin{align*}
  A &= \textstyle\sum_{\text{old}} D_{\F_{p-1}}(f)^2 \cdot
       (1/n^2 - 1/{n'}^2) && \text{(dilution --- always positive),}\\
  B &= (2/{n'}^2)\textstyle\sum D_{\F_{p-1}}(f) \cdot \delta(f)
       && \text{(cross term),}\\
  C &= (1/{n'}^2)\textstyle\sum \delta(f)^2
       && \text{(shift squared --- always positive),}\\
  \mathcal{D} &= (1/{n'}^2)\textstyle\sum_{\text{new}} D_{\F_p}(k/p)^2
       && \text{(new-fraction discrepancy --- always positive).}
\end{align*}
We use $\mathcal{D}$ (calligraphic) for the decomposition term to
distinguish it from the rank discrepancy function~$D(f)$.
Thus $\Delta W(p) \le 0$ (wobble worsens) if and only if
$B + C + \mathcal{D} \ge A$.

\begin{remark}[Structure of the dilution term]
The term $A = W(p-1) \cdot (n'^2 - n^2)/n^2$ is not an
independent quantity: it equals $W(p-1)$ scaled by a purely
geometric factor.  For large~$p$, since $n'/n = |\F_p|/|\F_{p-1}|
\approx p/(p-1)$, we have $n'^2/n^2 - 1 \approx 2/p$, so
$A \approx 2W(p-1)/p$.  The key theorem~$\mathcal{D}/A =
1 + O(1/p)$ therefore says: the new-fraction discrepancy
nearly cancels the dilution.  Substituting into the
decomposition, the two $W(p-1)$-dependent terms almost cancel:
\[
  n'^2\,\Delta W(p) \approx -(B + C).
\]
The net per-step change is thus determined by $B$ and $C$
alone---sums over \emph{existing} fractions' shifts---with
no dependence on the global state $W(p-1)$.  One could state
the main result directly as $n'^2\,\Delta W(p) \approx -(B+C)$
without naming $A$; the term~$A$ is retained because the
theorem $|\mathcal{D}/A - 1| = O(1/p)$ is its cleanest statement.
\end{remark}

\begin{remark}[Two-layer structure of the decomposition]
The terms $A$, $B$, $C$ in the decomposition
$\Delta W(p) = A - B - C - \mathcal{D}$ are defined
entirely from the existing fractions $\F_{p-1}$: they
involve only the current displacements $D_{\F_{p-1}}(f)$
and the shifts $\delta(f) = f - \{pf\}$.  No primality
of~$p$ is required for these terms individually.  The
connection to the arithmetic of the Riemann zeros enters
exclusively through~$\mathcal{D}$: the new fractions
$\{k/p : 1 \le k < p\}$ are equispaced precisely because
$\varphi(p) = p - 1$ (primality), and their squared
displacements in $\F_p$ carry the spectroscopic signal.
This separation---Farey-intrinsic identity plus arithmetic
input from primality---is what allows the per-step view
to go beyond a restatement of the primes-to-zeros chain.
\end{remark}

Computation through $p = 100{,}000$ reveals a striking near-cancellation:

\begin{observation}[Near-cancellation]
For all tested primes with $M(p) \le -3$:
\begin{enumerate}[(i)]
  \item $\mathcal{D}/A \in [0.97,\, 1.12]$: the new-fraction discrepancy cancels the dilution
    to within $O(1/p)$: specifically $\mathcal{D}/A = 1 + O(1/p)$.
  \item $B$ is \textbf{positive} for all tested $M(p) \le -3$ primes:
    the rank discrepancy $D$ and
    the shift $\delta$ are positively correlated in this regime.
  \item $C/A \in [0.05, 0.18]$ for the tested range.
  \item The combined ratio $(B + C + \mathcal{D})/A$ grows from ${\sim}1.4$ at
    $M = -3$ to ${\sim}3.0$ at $M = -14$.
\end{enumerate}
\end{observation}

\noindent A direct proof that $B \ge 0$ remains an open problem.
However, the Sign Pattern (Observation~\ref{thm:sign}) bypasses
this entirely: the proof shows $C + \mathcal{D} > A$
without requiring $B$'s sign, using a Cauchy--Schwarz bound on
the $\mathcal{D}/A$ ratio
combined with a rearrangement-inequality lower bound on~$C$.

\begin{remark}[Per-denominator approach fails]

A natural proof attempt is to decompose the correlation ratio
$R = \sum D \cdot \delta / \sum \delta^2$ by individual
denominators, writing $R = \sum_b R_b \cdot w_b$ with
weights $w_b = \delta_b^2 / \sum \delta^2$.  If each
$R_b > -1/2$, then $R > -1/2$ follows.  However, this
approach \emph{fails}: individual $R_b$ values can be
far below $-1/2$ (empirically as low as $-3$ for specific
denominator classes).  The bound $R > -1/2$ holds only
after cross-denominator cancellation aggregates the
contributions.
\end{remark}

\begin{remark}[Cross-denominator cancellation structure]
\label{rem:R-cancel}
Computation reveals that multiplication by~$p$ creates
\emph{more} cancellation in $R = \sum D \cdot \delta / \sum \delta^2$
than random permutations.  Specifically, $|R|$ uses only
${\sim}10\%$ of its Cauchy--Schwarz budget through a two-stage
mechanism: within-denominator cancellation (${\sim}63\%$)
and cross-denominator cancellation (${\sim}30\%$).  In tests
against 500 random permutations of the shift vector, the
actual $|R|$ falls \emph{below all random trials} ($z$-scores
from $-4$ to $-24$).  Structurally, $D$ is a global quantity
(rank among all fractions) while $\delta$ is local (determined
by $pa \bmod b$ within a single denominator), creating a
Weyl-sum--like pairing of smooth against scrambled.  When
$|M(p)|$ is large, within-denominator correlations become
positive, breaking the cancellation and explaining why
$|R|$ tracks $|M(p)|$.  Formalizing this cancellation
structure---perhaps via character sum techniques and the
Weil bound for Dedekind sums---is a candidate approach.
The reformulation $B + C = (-2/{n'}^2)\sum R(f)\,\delta(f)$
may connect to the theory of averaged Dedekind sums over
Farey fractions developed by
Minelli--Sourmelidis--Technau~\cite{Minelli2023a,Minelli2023b},
whose results on restricted averages over $(0,1/2)$ are
directly analogous to our per-denominator decomposition.

The Permutation Square-Sum Identity
(Lemma~\ref{thm:perm-sq}) makes this connection
precise: the reformulation $B + C = (-2/{n'}^2)\sum R(f)\,\delta(f)$
reduces the cross term to an inner product of
M\"obius-weighted fractional-part sums against
multiplicative displacements, which for fixed denominator~$b$
decomposes into exponential sums amenable to the
Weil bound (though the direct Kloosterman formulation
is invalid for this problem; see Conjecture~\ref{thm:shift-sq-asymp}).
\end{remark}

\subsection{Rubinstein--Sarnak connection}

The sigmoid dependence on $M(p)/\sqrt{p}$ (Section~\ref{sec:phenomenon})
means the overall rate of positive $\Delta W$ is governed by the distribution
of $M(p)/\sqrt{p}$ across primes.  This distribution depends on
the zeros of $\zeta(s)$ via the explicit formula for the Mertens
function, connecting our violation phenomenon to the
Rubinstein--Sarnak theory of Chebyshev's
bias~\cite{RubinsteinSarnak1994}.  In that framework, the
logarithmic density of integers where a prime race is ``won'' by
one residue class is determined by the zeros of Dirichlet
$L$-functions.  Analogously, the density of primes that
improve Farey uniformity may be expressible through the zeros
of~$\zeta(s)$---a connection that would link the per-step
geometry of Farey sequences to the Riemann Hypothesis in a
new way.





\section{Formal Verification}\label{sec:formal}
%% ================================================================

Key supporting results have been formally verified in Lean~4
using the Mathlib library, with the Aristotle automated theorem
prover~\cite{Aristotle2025} generating proof strategies.
The formalization spans 31~Lean~4 files containing 441~named
results (theorems, lemmas, definitions, and instances).
The core identity files---Bridge Identity, Injection Principle,
Displacement-Shift, Strict Positivity, and four-term
decomposition---compile with zero~\texttt{sorry} and no axiom
violations.  Two files contain~\texttt{sorry} in secondary
results: \texttt{MediantMinimality.lean} (the classical Farey
gap upper bound, used only for background) and
\texttt{SignTheorem.lean} (the universal sign conjecture,
which is retained as a placeholder for the disproved statement
at $p = 243{,}799$; the bounded-range version
$p \le 140$ is fully proved via \texttt{native\_decide}).  The files are:
\begin{itemize}
  \item \texttt{PrimeCircle.lean} (14 results): Farey cardinality,
    Ramanujan sums, wobble decomposition.
  \item \texttt{BridgeIdentity.lean} (24 results): Mertens function,
    coprime permutation, bridge identity proof chain.
  \item \texttt{CharacterBridge.lean} (9 results): character-weighted
    bridge identity for all Dirichlet characters.
  \item \texttt{InjectionPrinciple.lean} (10 results): formal proof
    of the Injection Principle (the Injection Principle).
  \item \texttt{GeneralInjection.lean} (11 results): injection
    principle for all~$N \ge 2$ (Proposition~the Generalized Injection Principle).
  \item \texttt{DisplacementShift.lean} (15 results):
    displacement-shift identity
    (Proposition~\ref{prop:disp-shift}).
  \item \texttt{DeltaCosine.lean} (6 results):
    $\delta$-cosine framework.
  \item \texttt{DenominatorSum.lean} (15 results):
    $\sum D(a/b) = -\varphi(b)/2$ per denominator.
  \item \texttt{StrictPositivity.lean} (15 results):
    strict positivity of $\sum \delta^2$
    via the rearrangement inequality.
  \item \texttt{CrossTermPositive.lean} (51 results):
    $B + C$ positivity framework, correlation ratio bounds,
    and four-term decomposition.
  \item \texttt{SubGapPermutation.lean} (12 results):
    sub-gap permutation identities for prime insertion.

  \item \texttt{SignTheorem.lean} (32 results,
    zero~\texttt{sorry}): Sign Pattern for all primes
    $13 \le p \le 113$ with $M(p) \le -3$ (verified by \texttt{native\_decide}), wobble
    monotonicity equivalence, correlation ratio $R(p) > -1/2$
    verified for all primes $11 \le p \le 83$,
    the structural theorem \texttt{bPlusC\_pos\_of\_corrRatio},
    exact values $R(11) = -1155/5974$ and $R(13) = 813/15872$,
    corrected ratio test (the original had a counterexample at $p = 7$),
    and the \texttt{dispRoughSquaredSum} definition showing
    Cauchy--Schwarz is too loose ($16{-}44\times$ gap (depending on~$p$; see Remark~\ref{rem:cs-gap})).
    The disproved universal conjecture has been commented out.
  \item \texttt{MediantMinimality.lean} (8 results):
    mediant minimality (mediant minimality),
    Farey gap formula (the classical Farey gap formula),
    and gap bounds for~$\F_N$.
  \item \texttt{MertensGrowth.lean} (8 results):
    computational witnesses for $|M(N)| > \sqrt{N}/2$
    and composite healing bounds.
  \item \texttt{FourierModeExploration.lean} (28 results):
    Fourier mode analysis of the cross term, per-denominator
    decomposition, and vanishing-denominator identities.
\end{itemize}
The Aristotle theorem prover~\cite{Aristotle2025}
autonomously constructed the bridge identity proof chain
(coprime residue permutation $\to$ Ramanujan sum at coprime
argument $\to$ bridge identity), the character-weighted
bridge, the Injection Principle, the displacement-shift
identity, and the mediant minimality theorem.
The bridge is additionally verified by
\texttt{native\_decide} for all primes $p \le 13$.

\subsection{Key verified results}

All of the following are fully proved with zero \texttt{sorry}:
\begin{enumerate}
  \item \textbf{Bridge Identity} (Theorem~\ref{thm:bridge}):
    $\sum_{f \in \F_{p-1}} e^{2\pi ipf} = M(p) + 2$.
    The complete proof chain---coprime residue permutation,
    Ramanujan sum at coprime argument, and the bridge
    itself---was proved autonomously by Aristotle.
  \item \textbf{Farey cardinality recurrence}:
    $|\F_N| = |\F_{N-1}| + \varphi(N)$.
  \item \textbf{Ramanujan sum}: $c_q(1) = \mu(q)$ for $q \ge 1$,
    generalized to $c_q(m) = \mu(q)$ when $\gcd(m,q) = 1$.
  \item \textbf{Coprime residue permutation}:
    multiplication by $p$ permutes $\{a : \gcd(a,b) = 1\}$



\smallskip\noindent
\textbf{Automation and reproducibility.}\quad
Proof chains were constructed autonomously by the
Aristotle theorem prover~\cite{Aristotle2025} and verified
by the Lean~4 kernel.
All source code (C, Python, Lean~4), computational data,
and figures are available at
\url{https://github.com/SaarShai/Primes-Equispaced}.

\section{Analytical Lower Bounds}\label{sec:lower-bounds}
%% ================================================================

This section collects four analytical results that strengthen
the Sign Pattern heuristic: a deficit minimality theorem for the
multiply-by-$2$ permutation, a spectral positivity result
connecting the Dedekind kernel to $L$-functions, an
asymptotic for the total shift-squared sum (conditional on a Dedekind sum estimate)
$\sum\delta_b^2$, and unconditional lower bounds on the
weighted discrepancy~$C_W(N)$.

\end{enumerate}

\subsection{Deficit minimality: $D_q(2) = q(q^2-1)/24$ is minimal}

For a prime~$q$ and multiplier $r \in \{2, \ldots, q-1\}$,
define the \emph{multiplicative deficit}
$D_q(r) = \sum_{a=1}^{q-1} a^2 - \sum_{a=1}^{q-1}
a \cdot (ra \bmod q)$.

\begin{theorem}[Deficit Minimality; classical, cf.\ Dedekind sum extrema]
For every prime $q \ge 3$ and every $r$ with $2 \le r \le q-1$:
\begin{equation}
  D_q(r) \;\ge\; D_q(2) \;=\; \frac{q(q^2-1)}{24},
\end{equation}
with equality if and only if $r = 2$ or
$r = (q+1)/2 = 2^{-1} \bmod q$.
\end{theorem}

\begin{proof}
Reduce to a Dedekind sum inequality: $D_q(r) \ge D_q(2)$
is equivalent to $s(r,q) \le s(2,q) = (q-1)(q-5)/(24q)$,
where $s(r,q)$ is the Dedekind sum.
For large multipliers $r > (q-1)/2$: the antisymmetry
$s(r,q) = -s(q-r,q)$ gives $s(r,q) \le 0 < s(2,q)$.
For small multipliers $2 \le r \le (q-1)/2$: apply Dedekind
reciprocity $s(r,q) + s(q,r) = (r^2+q^2+1)/(12rq) - 1/4$
and the permutation bound $s(a,b) \le s(1,b) = (b-1)(b-2)/(12b)$.
The resulting inequality reduces to
$2(r - 2)(r - (q+1)/2) \le 0$
for $2 \le r \le (q-1)/2$, with equality at the roots
$r = 2$ and $r = 2^{-1} \bmod q$.
Verified computationally for all primes $q \le 2{,}000$
(302~primes, ${\sim}300{,}000$ pairs, zero violations).
\end{proof}

\begin{remark}
The deficit $D_q(2)$ satisfies $D_q(2)/q = (q^2-1)/24$ for
prime~$q$.  Note: this deficit ratio is a component of the
four-term decomposition, not the full wobble change $\Delta W(q)$.  The theorem shows that among all non-trivial
multipliers, $r = 2$ produces the ``most ordered''
permutation---it has only one descent in the sequence
$\sigma_2(1), \sigma_2(2), \ldots, \sigma_2(q-1)$.
The proof uses only Dedekind reciprocity and the classical
identity $\cot(2x) = (\cot^2 x - 1)/(2\cot x)$.
\end{remark}

\subsection{Spectral positivity of the Dedekind kernel}

Define the Dedekind kernel $K_p(r) = s(r,p)$ and its
character transform
$\hat{K}_p(\chi) = \sum_{r=1}^{p-1} s(r,p)\,\chi(r)$
for Dirichlet characters $\chi \bmod p$.

\begin{proposition}[Spectral Factorization]
For every character $\chi \bmod p$,
\[
  \hat{K}_p(\chi) = B_{1,p}(\chi)\,B_{1,p}(\bar{\chi}),
\]
where $B_{1,p}(\chi) = \sum_{a=1}^{p-1}\chi(a)\,\bigl(\!\bigl(\tfrac{a}{p}\bigr)\!\bigr)$
and $((x)) = \{x\} - 1/2$ for $x \notin \Z$, and $((x)) = 0$ for $x \in \Z$.
\end{proposition}

\begin{proof}
Expand $\hat{K}_p(\chi) = \sum_r \chi(r) \sum_a ((a/p))\,((ra/p))$.
For fixed~$a$, substitute $m = ra$ and use $\chi(r) = \bar{\chi}(a)\chi(m)$.
The inner sum becomes $\bar{\chi}(a)\,B_{1,p}(\chi)$; summing over~$a$
gives $B_{1,p}(\chi)\,B_{1,p}(\bar{\chi})$.
\end{proof}

\begin{corollary}[Spectral Positivity]\label{cor:spectral-pos}
For every Dirichlet character $\chi \bmod p$:
\begin{enumerate}[(i)]
  \item If $\chi$ is even ($\chi(-1) = +1$), then $\hat{K}_p(\chi) = 0$.
  \item If $\chi$ is odd ($\chi(-1) = -1$), then
    $\hat{K}_p(\chi) = |B_{1,p}(\chi)|^2 \ge 0$.
  \item For every nonprincipal odd~$\chi$:
    \begin{equation}
      \hat{K}_p(\chi) = \frac{p}{\pi^2}\,|L(1,\chi)|^2.
    \end{equation}
\end{enumerate}
\end{corollary}

\begin{proof}
(i)~For even~$\chi$, pairing $a$ with $p-a$ gives $B_{1,p}(\chi) = 0$
since $((1-a/p)) = -((a/p))$ and $\chi(p-a) = \chi(a)$.
(ii)~For odd~$\chi$, $B_{1,p}(\bar\chi) = \overline{B_{1,p}(\chi)}$.
(iii)~For primitive odd~$\chi$, the functional equation gives
$L(1,\chi) = (\pi i/p)\,\tau(\chi)\,B_{1,p}(\bar\chi)$
with $|\tau(\chi)|^2 = p$, so
$|B_{1,p}(\chi)|^2 = (p/\pi^2)|L(1,\chi)|^2$.
Numerically verified for $p = 11, 13, 17, 23$.
\end{proof}

\begin{remark}
This spectral structure---exact vanishing on even characters,
positive squares on odd characters, and the explicit $L$-function
formula---provides a direct bridge from the Dedekind kernel
(which governs our per-denominator deficits) to analytic number
theory.  The total spectral mass satisfies
$\sum_{\chi\,\mathrm{odd}} \hat{K}_p(\chi) = (p-1)^2(p-2)/(12p)$,
giving an average odd eigenvalue of $(p-1)(p-2)/(6p) \sim p/6$.
\end{remark}

\subsection{Total shift-squared sum: $\sum\delta_b^2 = N^2/(2\pi^2) + O(N\log N)$}

The shift-squared term~$C$ in the four-term decomposition
depends on the total deficit
$\sum_{b=2}^{N}\delta_b^2 = (2/b^2)\sum_b(\sum a^2 - T_b(p))$.

\begin{conjecture}[Total Shift-Squared Asymptotic (conditional on Dedekind sum estimate)]\label{thm:shift-sq-asymp}
For every prime $p \ge 3$ with $N = p-1$:
\begin{equation}
  \sum_{b=2}^{N} \delta_b^2
  \;=\; \frac{N^2}{2\pi^2} + O(N\log N).
\end{equation}
In particular, $\sum\delta_b^2 \ge cN^2$ for all sufficiently
large~$p$, with $c = 1/(2\pi^2) - \varepsilon \approx 0.0507$.
\end{conjecture}

\begin{proof}[Heuristic derivation]
Write $\sum\delta_b^2 = 2(A(N) - B(N,p))$ where
$A(N) = \sum_b S_2(b)/b^2$ and $B(N,p) = \sum_b T_b(p)/b^2$.
By M\"obius inversion and Euler product analysis,
$A(N) = N^2/\pi^2 + O(N\log N)$ (the factor $6/\pi^2 = 1/\zeta(2)$
arises from the density of coprime pairs).
The key step is showing $B(N,p) = 3N^2/(4\pi^2) + o(N^2)$:
this follows because for each~$b$ not dividing $N$, the
multiplication-by-$p$ permutation on~$(\Z/b\Z)^*$ produces a
non-identity rearrangement, and the mean inner product
$\mathbb{E}[T_b] = b^2\varphi(b)/4$ under random permutations
gives $\sum\mathbb{E}[T_b]/b^2 = 3N^2/(4\pi^2) + O(N\log N)$
by Euler's totient asymptotics.
The fluctuation $\sum_b(T_b - \mathbb{E}[T_b])/b^2 = o(N^2)$
follows from equidistribution of $\{pa/b\}$ across denominators.

\emph{Analytic gap:} The fluctuation
$T_b(p) - \mathbb{E}[T_b]$ is not a Kloosterman sum
but a Dedekind-sum convolution:
$T_b(p) - \mathbb{E}[T_b] = b^2 \sum_{c|b} \mu(b/c)\,s(p,c)$.
For $b = p{-}1$, the fluctuation is $\Theta(b^3)$,
so a per-denominator bound of $O(b^{3/2+\varepsilon})$
is \emph{false}.  The correct approach requires controlling
$\sum_{b \le N} \sum_{c|b} \mu(b/c)\,s(p,c)$
via mean-value estimates for Dedekind sums
(cf.\ Conrey et al.~\cite{ConreyEtAl1996}).
This is why the result is stated as a conjecture.
Numerical confirmation: $\sum\delta_b^2/N^2$ converges to
$1/(2\pi^2) \approx 0.0507$ for $p = 97, 397, 997$
(values $0.0490, 0.0509, 0.0502$).
\end{proof}

\begin{remark}
The predicted leading constant is $1/(2\pi^2)$ (not just a lower bound),
arising from
$A(N) = N^2/\pi^2$ (coprime residue squares via M\"obius) and
$B(N) \sim 3N^2/(4\pi^2)$ (random permutation baseline).
Composites contribute $75.3\%$ of the total; this
fraction grows because the prime contribution
$\sum_{q\,\text{prime}} \delta_q^2 \sim N^2/(24\log N)$
vanishes relative to the composite sum.
\end{remark}

\subsection{Unconditional lower bounds on $C_W(N)$}

Define the weighted Farey discrepancy
$C_W(N) = (1/n)\sum_{f \in \F_N} D(f)^2$ where $n = |\F_N|$.

\begin{remark}[Elementary bound]
A simple Cauchy--Schwarz argument gives $C_W(N) \ge 1/4$ for all $N$,
but the following theorem provides a strictly stronger linear bound.
\end{remark}

\begin{theorem}[Linear Growth Lower Bound]\label{thm:cw-linear}
For all $N \ge 10$:
\begin{equation}
  C_W(N) \;\ge\; \frac{N}{28}.
\end{equation}
\end{theorem}

\begin{proof}
For denominators $b$ with $N/2 < b \le N$, the fraction
$1/b$ has exact rank $N - b + 1$ in~$\F_N$ (since the only
fractions less than~$1/b$ are $0/1$ and $1/b'$ with $b' > b$).
The displacement $|D(1/b)| = n/b - (N-b+1)$ satisfies
$|D(1/b)| \ge 0.05\,N$ for $b \in (3N/4, N]$,
using $n \ge 0.3\,N^2$ (from $|\F_N| = 3N^2/\pi^2 + O(N\log N)$).
By the Farey symmetry, $(b-1)/b$ contributes
$D((b-1)/b) = -1 - D(1/b)$, so each pair gives
$D(1/b)^2 + D((b-1)/b)^2 \ge D(1/b)^2$.

By the Cauchy--Schwarz inequality
($\sum x_i^2 \ge (\sum x_i)^2 / M$ for $M$~terms), applying Cauchy--Schwarz over all $M = \lfloor N/2 \rfloor$
terms in the range $b \in (N/2, N]$ and using the harmonic
sum $\sum_{b=\lceil N/2\rceil+1}^{N}1/b \ge \ln 2 - 2/N$:
$\sum D^2 \ge 0.01\,N^3$, hence
$C_W \ge 0.01\,N^3/(0.36\,N^2) > N/36$.
Including both $1/b$ and $(b-1)/b$ terms improves this to
$C_W \ge N/28$.  Verified for all $N \in [10, 2000]$
(minimum ratio $C_W/(N/28) = 1.88$ at $N = 10$).
\end{proof}

\begin{corollary}[Unconditional $p^2$ lower bound]
For every prime $p \ge 11$:
\[
  \sum_{f \in \F_p} E(f)^2 \;\ge\; \frac{p^2}{28},
\]
where $E(f) = \rank(f) - |\F_p|\cdot f$ is the Farey
discrepancy at the new fractions.  This gives
$C + \mathcal{D} \ge c\,p^2/{n'}^2$ for an explicit $c > 0$,
which is the key ingredient for closing the analytical tail
in the Sign Pattern.
\end{corollary}

\begin{remark}[Improved bounds via analytical rank formulas]
The constant $1/28$ in Theorem~\ref{thm:cw-linear} can likely be improved
using the analytical rank formulas of Tom\'as~Garc\'ia~\cite{TomasGarcia2022, TomasGarcia2024, TomasGarcia2025},
which provide exact expressions for the rank and local discrepancy of Farey
fractions with tighter estimates than the unit-fraction approach used here.
Under the Riemann Hypothesis, $C_W(N) \sim (1/2\pi^2)\,N\log N$,
so our unconditional bound $C_W \ge N/28$ is far from optimal.
\end{remark}


%% ================================================================
\section{The Sign Pattern, its Limits, and the Chebyshev Bias}\label{sec:sign-theorem}
%% ================================================================

The computational evidence reveals a striking pattern.

\begin{observation}[Computational Sign Pattern]\label{thm:sign}
For every prime $11 \le p \le 100{,}000$ with $M(p) \le -3$:
$\Delta W(p) < 0$.
\end{observation}

\noindent This is a computational result, verified by direct
computation in C (double precision, with tightest margin
$|\Delta W| = 7 \times 10^{-11}$ at $p = 92{,}177$) for all
$4{,}617$ qualifying primes in range.

\begin{remark}[The pattern fails at large $p$]\label{rem:failure}
Observation~\ref{thm:sign} does \emph{not} extend to all qualifying
primes.  We have verified $\Delta W(243{,}799) > 0$ (where
$M(243{,}799) = -3$), which disproves any universal sign claim
for $M(p) \le -3$ primes.  The correlation ratio
$R(p) = \sum D \cdot \delta / \sum \delta^2$ satisfies
$R(243{,}799) = -3.052$ (negative), driving $B + C < 0$ and
hence $\Delta W > 0$.
  This failure is not accidental: the Chebyshev phase argument
  $\gamma_1 \log(243{,}799) + \varphi \approx 180.6$~rad satisfies
  $\cos(\gamma_1 \log p + \varphi) \approx -0.033$, placing
  $p = 243{,}799$ precisely at a zero-crossing of the leading
  oscillation.  The sign flip is an expected consequence (given the empirical phase correlation) of
  the strongly correlated wave crossing the axis.  For all $4{,}617$ primes in
Observation~\ref{thm:sign}, $R > 0$ (minimum $R = 0.0068$ at $p = 64{,}781$).

The question of \emph{why} $R$ turns negative at $p = 243{,}799$
is answered by the Chebyshev bias framework of
Section~\ref{sec:sign-theorem}: the sign of $\Delta W(p)$ is
strongly correlated to $\cos(\gamma_1 \log p + \varphi)$, so both
both signs occur infinitely often (unconditional, Ingham's theorem~\cite{Ingham1942}).
Under GRH~+~LI, the leading oscillation term is equidistributed,
suggesting the density of negative steps approaches~$1/2$; however,
the exact limiting density may differ due to lower-order bias terms
(cf.\ Rubinstein--Sarnak~\cite{RubinsteinSarnak1994}).
\end{remark}

\begin{figure}[t]
\centering
\includegraphics[width=\textwidth]{phase_lock_visualization.png}
\caption{Phase-lock: the smoothed correlation ratio $R(p)$
tracks the leading zeta-zero oscillation
$\cos(\gamma_1\log p + \varphi)$.}
\label{fig:phase correlation}
\end{figure}

A stronger computational result holds for small primes:

\begin{theorem}[Sign Pattern --- small primes]\label{thm:sign-small}
For every prime $11 \le p \le 113$:
$\Delta W(p) < 0$,
regardless of the value of~$M(p)$.
\end{theorem}

\noindent This is formally verified in Lean~4
(\texttt{sign\_theorem\_all\_le\_113}) by
\texttt{native\_decide} on exact rational computations.
The small primes $p = 2, 3, 5, 7$ all have $\Delta W(p) > 0$;
the threshold for universal $\Delta W < 0$ is therefore
$p \ge 11$, not $p \ge 5$ as previously conjectured.
The result is unconditional (no assumption on~$M(p)$) and
establishes that the $M(p) \le -3$ condition in
Observation~\ref{thm:sign} is needed only for $p > 113$.

The proof relies on two analytical ingredients.

\begin{proposition}[$\mathcal{D}/A$ ratio]\label{prop:DA-ratio}
For every prime $p \ge 3$ with $M(p) \le -3$,
\[
  \frac{\mathcal{D}}{A} = 1 + O(1/p),
\]
where the implicit constant depends only on $|M(p)|/p$.
This follows from wobble conservation: the total wobble
$\sum_{f \in \F_{p-1}} w(f) = 0$ forces the new-fraction
discrepancy sum to nearly cancel the dilution.
\end{proposition}

\begin{remark}[Quadratic disruption and the prime number theorem]

Each prime~$p$ inserts $\varphi(p) = p{-}1$ fractions, creating
disruption proportional to~$p^2$ (from Proposition~\ref{prop:D1p}).
Primes appear with density $\sim 1/\log p$, so the total disruption
per unit interval is $p^2/\log p$, which grows.
Primes become \emph{rarer} but individually
\emph{more disruptive} --- this balance is the Farey-sequence
analogue of the prime number theorem.
\end{remark}

\begin{proposition}[Strict positivity of $C$]\label{prop:strict-pos}
For every prime $p \ge 5$, the shift-squared term satisfies $C > 0$.
Equivalently, for each denominator $b < p$,
$T_b(p) \le \sum_{a} a^2$ with equality if and only if
$p \equiv 1 \pmod{b}$, and the latter cannot hold simultaneously
for all $b$, so the strict inequality $C > 0$ follows by the
rearrangement inequality.
\end{proposition}

\begin{proof}[Proof of Observation~\ref{thm:sign}]
Direct computation in~C (double precision) verifies
$\Delta W(p) < 0$ for all $4{,}617$ primes in the stated
range.  The tightest margin is $|\Delta W| = 7 \times 10^{-11}$
at $p = 92{,}177$, four orders of magnitude above float64
rounding.  For $p \le 200$, exact rational arithmetic
independently confirms the sign.  The following analytical
ingredients support the geometry of why the sign is almost
always negative.

\begin{enumerate}[(i)]
\item \emph{$\mathcal{D}/A = 1 + O(1/p)$}
(Proposition~\ref{prop:DA-ratio}): the new-fraction discrepancy
nearly cancels the dilution, reducing the sign question to
whether $B + C = 2\sum D \cdot \delta + \sum \delta^2 > 0$.

\item \emph{$C > 0$ strictly}
(Proposition~\ref{prop:strict-pos}, rearrangement inequality).

\item \emph{$C' \ge c\,N^2$}
By Conjecture~\ref{thm:shift-sq-asymp},
$\sum\delta_b^2 = N^2/(2\pi^2) + O(N\log N)$,
so $C'$ grows quadratically.  Empirically $C' \ge 0.035\,p^2$
for all primes $p \ge 11$.

\item Computational verification confirms
$\mathcal{D}/A + C/A > 1$ for all primes $p \le 3{,}000$
with $M(p) \le -3$ (minimum ratio $1.096$ at $p = 2{,}857$).
\end{enumerate}

\noindent Note: items~(i)--(iii) show that $C$ alone is
growing, but whether $B + C > 0$ (i.e., whether the cross
term $B = 2\sum D \cdot \delta$ is not too negative) is
not guaranteed analytically.  At $p = 243{,}799$, $B + C < 0$
(Remark~\ref{rem:failure}), so the extension beyond the
computational range requires the Chebyshev bias analysis of
Section~\ref{sec:sign-theorem}.
\end{proof}

\begin{remark}[The correlation ratio $R$ and open conjecture]
The sign of $\Delta W(p)$ is equivalent to the sign of
$B + C = 2\sum D \cdot \delta + \sum \delta^2 = \sum\delta^2 \cdot (1 + 2R)$,
where $R(p) = \sum D \cdot \delta \,/\, \sum \delta^2$ is
the \emph{correlation ratio}.  When $\mathcal{D}/A \approx 1$ (Proposition~\ref{prop:DA-ratio}),
the sign of $\Delta W(p)$ is approximately determined by
whether $R(p) > -1/2$.

\textbf{Computational findings:}
For all $4{,}617$ qualifying primes $p \le 100{,}000$,
$R(p) > 0$ (minimum $R = 0.0068$ at $p = 64{,}781$; $R$ grows
with $|M(p)|$, reaching $R > 73$ at $M(p) = -88$).
At $p = 243{,}799$ (the first known sign-flip),
$R(243{,}799) = -3.052$, confirming $B + C < 0$ there.

\begin{conjecture}[$R > 0$ for almost all qualifying primes]
For all but finitely many primes $p$ with $M(p) \le -3$,
$R(p) > 0$.
\end{conjecture}

\noindent This conjecture is open.  All analytical proof
paths attempted (Weil bounds, Abel summation,
Fourier analysis, induction) have failed
(see supplementary materials on GitHub).
The Chebyshev bias framework (Section~\ref{sec:sign-theorem})
explains \emph{why} $R$ can turn negative at specific primes
(such as $p = 243{,}799$) and predicts that sign flips occur.
\end{remark}

\begin{proposition}[Dominance of $1/p$]\label{prop:D1p}
For every prime $p\ge3$, the rank deviation of the fraction
$1/p$ in~$\F_p$ satisfies
\[
  D\!\left(\tfrac{1}{p}\right)
  = 1 - \frac{|\F_p|}{p}
  \;\sim\; -\frac{3p}{\pi^2}
  \qquad(p\to\infty).
\]
The product $D(1/p)\cdot(1/p) \to -3/\pi^2\approx -0.304$,
where $g(1/p) = 1/p$ denotes the gap from $0/1$ to $1/p$ in~$\F_p$.
This
is asymptotically a universal constant.  Computationally, this
single fraction accounts for $65{-}73\%$ of the total
$|\sum D\cdot\delta|$ for $p\le199$.
\end{proposition}

\begin{proof}
The fraction $1/p$ enters~$\F_p$ at zero-based index $j=1$
(only $0/1$ precedes it, since $1/(p{-}1) > 1/p$ for $p\ge3$).
With $n = |\F_p|$, the rank discrepancy is
$D(1/p) = 1 - |\F_p|/p$.
The asymptotic follows from Mertens' estimate
$\sum_{k=1}^{N}\varphi(k) = 3N^2/\pi^2 + O(N\log N)$~\cite{HardyWright2008}.
\end{proof}

\begin{observation}[Damage--response decomposition]
Define the \emph{damage ratio}~$R_{\mathrm{new}}(p)$ by summing
$D\cdot g$ (where $g(f) = f - f_{\mathrm{prev}}$ is the Farey gap) over new fractions entering~$\F_p$
(where $\delta_1 = f - f_{\text{prev}}$ in $\F_p$), and the
\emph{response ratio}~$R(p)$ by summing $D\cdot\delta_2$ over
existing fractions in~$\F_{p-1}$ ((where $\delta(a/b) = (a - pa\bmod b)/b$ is the insertion
deviation of Definition~2.2)).

For all $4{,}617$ primes with $M(p)\le-3$ and $p\le100{,}000$:
$R_{\mathrm{new}}(p) < 0$ always (new fractions damage regularity) and
$R(p) > 0$ always (existing fractions overcompensate).
The net $\Delta W(p) < 0$ arises because the response dominates.

\smallskip\noindent
The sign alignment underlying $R>0$ is weak: at $p=31$,
only $59\%$ of old fractions have
$\sgn(D)\cdot\sgn(\delta_2)>0$.
Positivity is driven by magnitude: fractions with large
positive~$D$ also have large positive~$\delta_2$.
An analytical proof of $R>0$ remains open.
\end{observation}

\begin{figure}[ht]
\centering
\includegraphics[width=\textwidth]{fig_damage_response.png}
\caption{The damage/response mechanism.
\textbf{(a)}~Decomposition of $\Delta W$ for primes with
$M(p) \le -3$: new-fraction damage ($\mathcal{D}$, red) and
shift-squared ($C$, orange) push $\Delta W$ positive, while
the cross-term response ($B$, blue) partially compensates.
Green diamonds show net~$\Delta W$.  The darker stripe within
$\mathcal{D}$ marks the $D(1/p)^2$ contribution.
\textbf{(b)}~Share of $\mathcal{D}$ attributable to the single
fraction~$1/p$ (measured as $D(1/p)^2 / \mathcal{D}$);
the $D \cdot \delta$ product share is higher at $65$--$73\%$
(Proposition~\ref{prop:D1p}).}
\label{fig:damage-response}
\end{figure}




\noindent The conjecture is motivated by a density argument:
$\varphi(n)/n$ measures the fraction of integers $\le n$ coprime
to~$n$, hence the density of new fractions entering~$\F_n$.
When this density is below~$1/3$, the new fractions distribute
sufficiently evenly to strongly favor healing (though counterexamples exist
at $n = 438, 546, 570$).  The threshold~$1/3$
corresponds to $n$ having at least two prime factors below~$3$
(i.e., $6\mid n$), ensuring a rich divisibility structure.

%% ================================================================
\begin{remark}[Information compression]
The bridge identity compresses the geometric information of
$|\F_N| \sim 3N^2/\pi^2$ Farey fractions into a single integer
$M(N)$ via Ramanujan sum collapse --- a $2{,}807 : 1$ reduction
at $N=96$.  The mutual information between $M(p)$ and
$\sgn(\Delta W(p))$ is $\hat{I} \approx 0.79$~bits
(empirical, $p \le 100{,}000$), with
$\hat{P}(\Delta W < 0 \mid M < 0) = 99.98\%$.
The per-step entropy change $\Delta H$ correlates with
$\Delta W$ at $r = -0.914$ (exact, $p \le 200$).
\end{remark}


% [GK concentration + figure moved to supplementary]

\subsection{Connection to zeta zeros via the explicit formula}
\label{sec:explicit-connection}

The Bridge Identity (Theorem~\ref{thm:bridge}) connects the
cross-term~$B$ to the Mertens function:
$\sum_{f \in \F_{p-1}} e^{2\pi i p f} = M(p) + 2$.
Under GRH, the explicit formula for~$M(x)$
gives~\cite{Edwards1974}
\begin{equation}\label{eq:explicit}
  M(x) = \sum_{\rho} \frac{x^{\rho}}{\rho\,\zeta'(\rho)}
  - 2 + \text{(lower order)},
\end{equation}
where the sum runs over nontrivial zeros
$\rho = \tfrac{1}{2} + i\gamma$ of~$\zeta(s)$.
Writing $x^{\rho} = \sqrt{x}\,e^{i\gamma\log x}$,
the Mertens function at a prime~$p$ oscillates as
\[
  M(p) \;\approx\;
  2\sqrt{p}\sum_{k}\frac{\cos(\gamma_k\log p + \phi_k)}
  {|\rho_k\,\zeta'(\rho_k)|},
\]
where $\phi_k = -\arg(\rho_k\,\zeta'(\rho_k))$.
The dominant term ($k = 1$, $\gamma_1 \approx 14.13$) produces
a correlation $r = 0.77$ between $\sgn(\Delta W(p))$ and
$-\sgn(\cos(\gamma_1\log p + \varphi))$ across
$4{,}617$ qualifying primes
(Section~\ref{sec:sign-theorem}).

\begin{remark}[Conditional nature]
The derivation above is conditional on GRH (for the explicit
formula) and assumes term-by-term dominance by the leading zero.
An unconditional proof that $\Delta W(p)$ oscillates with zeta
zeros remains an open problem; the primary gap is an analytical
proof that $B > 0$ for all primes with $M(p) \le -3$
(Section~\ref{sec:phenomenon}).
\end{remark}


%% ================================================================
\section{The Farey Spectroscope}\label{sec:spectroscope}
%% ================================================================

The explicit formula connects $\Delta W(p)$ to the nontrivial zeros
of~$\zeta(s)$.  We exploit this in a new direction: rather than
\emph{assuming} the zeros to predict $\Delta W(p)$, we use Farey
discrepancy data to \emph{detect} the zeros.

\begin{definition}
The \emph{Farey spectral function} is
\[
  F(\gamma) \;=\; \Bigl|\,\sum_{\substack{p\;\text{prime}\\M(p)\le-3}}
  R(p)\,p^{-1/2-i\gamma}\Bigr|^2,
\]
where $R(p)$ is the insertion-deviation correlation ratio
(restricting to primes with $M(p)\le -3$ selects those where
the dilution term~$\mathcal{D}$ is small enough relative to~$A$
for the cross-term sign to dominate; the spectroscope also detects $\gamma_1$ without
this filter (using $M(p)/\sqrt{p}$ as weight on all primes;
see Figure~\ref{fig:null-control}), confirming the signal
is not an artifact of the subset selection)
The filter selects a dense subset of primes
($\approx 48\%$ of all primes $p \le 100{,}000$),
preserving the fundamental oscillation frequencies.
(the insertion-deviation ratio from Section~\ref{sec:defs}).
\end{definition}

\begin{figure}[t]
\centering
\includegraphics[width=\textwidth]{spectroscope_null_battery.png}
\caption{Null hypothesis battery ($500$ trials).
\textbf{Top:} real $F(\gamma)$ (red) towers above the
$95\%$ confidence envelope from shuffled-$R$ null spectra
(grey).  Dashed lines mark known zeta zeros.
\textbf{Bottom:} signal-to-null ratio exceeds $40\times$
at $\gamma_1$.  Formal tests: shuffled-$R$ $p < 0.002$
($z = 117.6$), Gaussian-$R$ $p < 0.002$ ($z = 106.5$),
wrong-sign primes produce zero signal,
zero false peaks above $50\%$ threshold.}
\label{fig:null-control}
\end{figure}

\begin{conjecture}[Spectral peaks, GRH]
Assume GRH with simple zeros.  Then $F(\gamma)$ has local maxima
at $\gamma = \gamma_k$ with peak heights satisfying
$F(\gamma_k) \propto |(\rho_k\,\zeta'(\rho_k))^{-1}|^2$
as the number of primes tends to infinity.
\end{conjecture}

\noindent\textbf{Heuristic.}\quad
Under GRH, the Farey counting error admits
$E(N) \sim \sum_\rho N^\rho/(\rho\,\zeta'(\rho))$~\cite{Edwards1974},
The substitution of $R(p)$ for the Mertens error is empirical:
our data shows $R(p)$ isolates the arithmetic ``damage'' of prime
insertion, which is proportional to the Mertens bias $M(p)/\sqrt{p}$.
We use $R(p)$ rather than the raw $\Delta W(p)$ because
$R(p)$ isolates the cross-term that carries the oscillatory
signal, stripping away the near-cancellation of $A$ and $\mathcal{D}$.
Because $R(p)$ strongly tracks $M(p)/\sqrt{p}$,
the spectroscope is mathematically akin to inverting the
explicit formula for~$M(x)$; its novelty lies in the fact
that this arithmetic signal is perfectly captured by the
purely geometric step-by-step discrepancy of Farey fractions.

so $R(p)\,p^{-1/2}\approx\sum_k c_k\,p^{i\gamma_k}$
with $c_k = (\rho_k\,\zeta'(\rho_k))^{-1}$.
At $\gamma = \gamma_j$ the $k{=}j$ term sums constructively
($\sum_p 1 \sim N/\log N$) while cross-terms
oscillate; a rigorous bound on the cross-term cancellation
would require pair-correlation estimates beyond what
is currently known unconditionally.
The heuristic is strongly supported by computation:
the predicted and observed amplitude ratios correlate at
$r = 0.997$ ($p$-value $< 10^{-8}$, $n = 10$ zeros), though the
monotonic decrease of both sequences inflates the
correlation (see~\S\ref{sec:open} for discussion;
after log-detrending, the residual correlation is
$r \approx -0.18$, confirming the raw~$r$ reflects shared
power-law decay).

\noindent\textbf{Detection of $\boldsymbol{\gamma_1}$.}\quad
With $N=3{,}829$ qualifying primes ($p\le83{,}773$),
$F(\gamma)$ detects the first three nontrivial zeros:
$\gamma_1\approx 14.08$ (error~$0.4\%$),
$\gamma_2\approx 20.86$ (error~$0.8\%$), and
$\gamma_3\approx 24.94$ (error~$0.3\%$)
(Figure~\ref{fig:spectroscope-preview}).

\smallskip\noindent\textbf{Amplitude matching.}\quad
The predicted ratios $|c_k/c_1|^2$ and observed
$F(\gamma_k)/F(\gamma_1)$ correlate at $r = 0.997$
across the first ten zeros
(Figure~\ref{fig:amplitude-matching}).

\smallskip\noindent\textbf{Extension to $\boldsymbol{L}$-functions.}\quad
Twisting by a Dirichlet character $\chi$,
$F_\chi(\gamma)=|\sum_p \chi(p)\,R(p)\,p^{-1/2-i\gamma}|^2$,
detects zeros of $L(s,\chi)$.  With $\chi_4$
(the non-principal character mod~$4$), peaks appear near
$\gamma'\approx 6.16$ and $16.13$, matching known zeros
of $L(s,\chi_4)$.  Since $\zeta_{\Q(i)}(s) = \zeta(s)\cdot L(s,\chi_4)$,
the untwisted and twisted spectroscopes together detect
zeros of the Dedekind zeta function of $\Q(i)$.
A systematic survey of all $108$ primitive characters
$\chi\!\pmod{q}$ for $q \le 20$ confirms that the
spectroscope detects peaks for every character tested
(see supplementary materials for the full survey).

\smallskip\noindent\textbf{Pair correlation.}\quad
The autocorrelation of $F(\gamma)$ detects four
pairwise zero differences (mean matching error~$0.15$),
each exceeding the $3\sigma$ null envelope from
$100$ shuffled-$R$ trials.  The strongest lag
$\approx 6.71$ matches $\gamma_2 - \gamma_1 = 6.89$,
connecting to the Montgomery pair correlation conjecture.

\smallskip\noindent\textbf{Convergence.}\quad
The spectroscope locates $\gamma_1$ to $0.34\%$
with $200$~qualifying primes.
Under GRH, the signal-to-noise ratio grows as $O(N)$
and the detection error as
$O(1/(\sqrt{N}\log p_{\max}))$, matching the Rayleigh
resolution limit for a periodogram on log-primes.
Unlike the Riemann--Siegel formula (which evaluates $\zeta(s)$
directly and converges exponentially faster with correction terms),
the spectroscope uses only Farey discrepancy data from primes.
The two methods are complementary: they operate on fundamentally
different inputs and provide complementary computation
(Figure~\ref{fig:juxtaposition}).

\begin{figure}[t]
\centering
\includegraphics[width=\textwidth]{farey_vs_classical_zeros.png}
\caption{Top: Hardy's $Z$-function, with zeros at sign changes.
Bottom: Farey spectral function $F(\gamma)$, with zeros as peaks.
Same zeros, different computational paths.}
\label{fig:juxtaposition}
\end{figure}

\begin{figure}[t]
\centering
\includegraphics[width=0.8\textwidth]{amplitude_matching_scatter.png}
\caption{Predicted $|c_k/c_1|^2$ versus observed $F(\gamma_k)/F(\gamma_1)$
for the first ten zeros.  Correlation $r = 0.997$.}
\label{fig:amplitude-matching}
\end{figure}



%% ================================================================

\begin{remark}[Status of main claims]
\begin{center}
\begin{tabular}{lll}
\toprule
Claim & Status & Method \\
\midrule
Bridge Identity & Proved & Algebraic + Lean~4 \\
Universal Exponential Sum & Proved & Algebraic + Lean~4 \\
Displacement-Cosine Identity & Proved & Bridge Identity + antisymmetry \\
Four-term decomposition & Proved & Algebraic + Lean~4 \\
$D(1/p) = 1 - |\F_p|/p$ & Proved & Elementary \\
$R(p) > 0$ for $p \le 106{,}649$ & Computational & C program \\
$\gamma_1,\gamma_2,\gamma_3$ detection & Computational & Spectroscope \\
Chebyshev phase correlation & Computational & $R = 0.77$ \\
Shift-squared asymptotic & Conjecture & Dedekind gap open \\
Spectral peaks $\propto |c_k|^2$ & Conjecture (GRH) & Heuristic \\
\bottomrule
\end{tabular}
\end{center}
\end{remark}

\section{Open Questions}\label{sec:open}

\begin{remark}[Structure of the Cauchy--Schwarz gap]\label{rem:cs-gap}
For $p = 997$ ($|\F_{996}| = 301{,}651$), the correlation ratio
$R(p) = 2.05$ while the Cauchy--Schwarz bound gives
$|R| \le 34.1$, a looseness factor of $16.6\times$.
Decomposing the cancellation:
approximately $75\%$ of the raw $|D\cdot\delta|$ mass
cancels \emph{within} fixed denominators (the products
$D(a/b)\cdot\delta(a/b)$ alternate wildly in sign as~$a$
varies over coprime residues), while only $25\%$ additional
cancellation occurs across denominators.

Among the $996$ active denominators, $832$ have positive
per-denominator sums and $152$ have negative sums,
so cross-denominator cancellation is secondary.
A tighter bound on~$R$ should therefore target the
within-denominator oscillation, potentially using
per-denominator Weil-type estimates for the correlation
between the Dedekind-sum structure of~$D$ and the
modular structure of~$\delta$.
\end{remark}

%% ================================================================

\begin{enumerate}
  \item \textbf{Threshold improvement.}
    The Sign Pattern (Observation~\ref{thm:sign}) shows computationally that
    $M(p) \le -3 \implies \Delta W(p) < 0$.
    Can the threshold be improved to $M_0 = -2$?
    The counterexample at $p = 92{,}173$ ($M = -2$) shows it
    cannot hold unconditionally, but a density-one result
    (holding for all but finitely many exceptions) may be achievable.

  \item \textbf{The cross-term sign.}
    The Sign Pattern bypasses the cross term $B = (2/n'^2)\sum D
    \cdot \delta$ entirely, proving $C + \mathcal{D} > A$ without
    $B$'s sign.  While $B > 0$ holds
    for all tested primes with $M(p) \le -3$,
    $B + C$ can be \emph{negative} for primes with large positive
    Mertens values (e.g.\ $B + C < 0$ at $p = 1399$, $M = +8$).
    The Permutation Square-Sum Identity
    (Lemma~\ref{thm:perm-sq}) reformulates the cross term as
    $B + C = (-2/{n'}^2)\sum R(f)\,\delta(f)$ where $R$ is a
    M\"obius-weighted fractional-part sum, connecting
    the sign question to Dedekind-sum estimates.

  \item \textbf{Other discrepancy measures.}  Does the per-step
    decomposition yield analogous phenomena for $L^1$ or $L^\infty$
    Farey discrepancy?

  \item \textbf{Violation density.}  The violation probability
    depends on $M(p)/\sqrt{p}$ via a sigmoid function.  Can this
    sigmoid be expressed analytically?  Does the overall violation
    density converge to a constant determined by the distribution
    of $M(p)/\sqrt{p}$?

  \item \textbf{Geometric growth bound (heuristic).}  If the
    empirical sigmoid persists, then both positive and negative
    $\Delta W$ occur infinitely often among primes, which would
    force $M$ to oscillate with amplitude at least
    $c\sqrt{p}$ for some $c > 0$.  This is a known result
    \cite{Ingham1942}, but a geometric proof via the sigmoid would
    be new.  Making this rigorous requires proving the
    sigmoid analytically, which reduces to the anticorrelation
    problem of question~2.

  
  \item \textbf{Growth rate of $N \cdot W(N)$.}
    The Franel--Landau theorem gives $W(N) = O(N^{-1+\varepsilon})$
    under RH, but the precise sub-polynomial rate is unknown.
    Computationally, $N \cdot W(N)$ grows from $\approx 0.28$ at
    $N = 2$ to $\approx 0.87$ at $N = 2000$; the ratio
    $N \cdot W(N) / \log N$ decreases from $0.119$ to $0.086$ over
    the same range, not converging to a constant.  What is the
    true asymptotic of $W(N)$---unconditionally and under RH?

  \item \textbf{Steinerberger optimality.}
    Steinerberger~\cite{Steinerberger2019} conjectures that greedily
    constructed sequences achieve discrepancy $O((\log N)/N)$,
    matching Schmidt's lower bound.  The Farey sequence is not
    greedy, yet the same rate $W(N) \sim C (\log N)/N$ is suggested
    by data (though not confirmed).  Is Farey discrepancy
    asymptotically Schmidt-optimal, despite being arithmetically
    (not greedily) defined?

% [Open Questions moved to supplementary]
\end{enumerate}


%% ================================================================
\section*{AI Use Declaration}

In accordance with STM Association recommendations (September 2025),
we declare the following uses of AI in preparing this manuscript:

\begin{itemize}
\item \textbf{Writing and drafting} (Activity~2): Anthropic's Claude
  (Opus~4.6, Sonnet~4.6) was used to draft and revise manuscript text,
  including theorem statements, proof sketches, and computational descriptions.
  All mathematical content was verified by the human author.

\item \textbf{Code generation and computation} (Activity~5): Claude,
  Google Gemma~4 (26B), and Alibaba Qwen~3.6 Plus were used to write
  Python and C programs for numerical experiments.  All code was reviewed
  and results independently verified.

\item \textbf{Data visualization} (Activity~6): Claude generated all
  12~figures from computational data.  Figures were reviewed for accuracy.

\item \textbf{Formal verification}: The Aristotle automated theorem
  prover (harmonic.fun) generated Lean~4 proofs.  OpenAI Codex performed
  adversarial code review of claimed results.

\item \textbf{Literature search} (Activity~7): Gemma~4 and Qwen~3.6
  assisted with identifying relevant references.
\end{itemize}

\noindent The human author (S.~Shai) takes full responsibility for the
correctness of all mathematical claims and the integrity of the work.

\section*{Acknowledgments}
%% ================================================================

The formal verification was carried out using the Aristotle
automated theorem prover~\cite{Aristotle2025}.  The compact
cross-term formula (Theorem~\ref{thm:cross-term}) was discovered
during a collaborative session with a GPT-based agent.
Computation was performed in C using the Farey next-term
generator, with results verified by Python exact rational
arithmetic for small primes.  We thank the Mathlib community
for maintaining the Lean~4 library that made the formalization
possible.


\section*{Data and Code Availability}

All source code (C, Python, Lean~4), computational data, and
figures are available at
\url{https://github.com/SaarShai/Primes-Equispaced}.

%% ================================================================
%% REFERENCES
%% ================================================================

\FloatBarrier
\begin{thebibliography}{99}

\bibitem{ConreyEtAl1996}
J.~B.~Conrey, E.~Fransen, R.~Klein, and C.~Scott,
\emph{Mean values of Dedekind sums},
J.\ Number Theory \textbf{56} (1996), 214--226.


\bibitem{Ingham1942}
A.~E. Ingham,
\emph{On two conjectures in the theory of numbers},
Amer.\ J.\ Math.\ \textbf{64} (1942), 313--319.

\bibitem{BCZ2000}
F.~P. Boca, C.~Cobeli, and A.~Zaharescu,
\emph{Distribution of lattice points visible from the origin},
Comm.\ Math.\ Phys.\ \textbf{213} (2000), 433--470.

\bibitem{Aristotle2025}
Harmonic,
\emph{Aristotle: Automated Theorem Proving for Lean~4},
\url{https://aristotle.harmonic.fun}, 2025.

\bibitem{BPS2015}
A.~I.~Bobenko, U.~Pinkall, and B.~A.~Springborn,
``Discrete conformal maps and ideal hyperbolic polyhedra,''
\emph{Geometry \& Topology}, vol.~19, pp.~2155--2215, 2015.
arXiv:1005.2698.

\bibitem{Franel1924}
J.~Franel,
``Les suites de Farey et le probl\`eme des nombres premiers,''
\emph{G\"ottinger Nachrichten}, pp.~198--201, 1924.

\bibitem{Landau1924}
E.~Landau,
``Bemerkungen zu der vorstehenden Abhandlung von Herrn Franel,''
\emph{G\"ottinger Nachrichten}, pp.~202--206, 1924.

\bibitem{RubinsteinSarnak1994}
M.~Rubinstein and P.~Sarnak,
``Chebyshev's bias,''
\emph{Experimental Mathematics}, vol.~3, no.~3, pp.~173--197, 1994.


\bibitem{Edwards1974}
H.~M.~Edwards,
\emph{Riemann's Zeta Function},
Academic Press, 1974.
Reprinted by Dover, 2001.

\bibitem{Ramanujan1918}
S.~Ramanujan,
``On certain trigonometrical sums and their applications in the
theory of numbers,''
\emph{Transactions of the Cambridge Philosophical Society},
vol.~22, pp.~259--276, 1918.

\bibitem{HardyWright2008}
G.~H.~Hardy and E.~M.~Wright,
\emph{An Introduction to the Theory of Numbers}, 6th ed.,
Oxford University Press, 2008.

\bibitem{Kanemitsu1982}
S.~Kanemitsu, R.~Sita~Rama~Chandra~Rao, and A.~Siva~Rama~Sarma,
``On some identities involving Ramanujan sums,''
\emph{The Mathematics Student}, vol.~50, pp.~110--123, 1982.

\bibitem{Duke1988}
W.~Duke,
``Hyperbolic distribution problems and half-integral weight Maass forms,''
\emph{Inventiones Mathematicae}, vol.~92, pp.~73--90, 1988.

\bibitem{Walfisz1963}
A.~Walfisz,
\emph{Weylsche Exponentialsummen in der neueren Zahlentheorie},
VEB Deutscher Verlag der Wissenschaften, Berlin, 1963.

\bibitem{KarvonenZhigljavsky2025}
T.~Karvonen and A.~Zhigljavsky,
``Maximum mean discrepancies of Farey sequences,''
arXiv:2407.10214, 2025.

\bibitem{TomasGarcia2025}
R.~Tom\'as~Garc\'ia,
``New analytical formulas for the rank of Farey fractions and
estimates of the local discrepancy,''
\emph{Mathematics}, vol.~13, no.~1, art.~140, 2025.

\bibitem{TomasGarcia2022}
R.~Tom\'as~Garc\'ia,
``Partial Franel sums,''
\emph{Journal of Integer Sequences}, vol.~25, art.~22.1.5, 2022.

\bibitem{TomasGarcia2024}
R.~Tom\'as~Garc\'ia,
``Improved rank formulas for Farey fractions,''
\emph{Integers}, vol.~24, art.~y63, 2024.

\bibitem{Steinerberger2019}
S.~Steinerberger,
``Dynamically defined sequences with small discrepancy,''
arXiv:1902.03269, 2019.

\bibitem{Minelli2023a}
P.~Minelli, A.~Sourmelidis, and N.~Technau,
``Bias in the number of steps in the Euclidean algorithm
and a conjecture of Ito on Dedekind sums,''
\emph{Mathematische Annalen} \textbf{387} (2023), 291--320.
arXiv:2206.03214.

\bibitem{Minelli2023b}
P.~Minelli, A.~Sourmelidis, and N.~Technau,
``On restricted averages of Dedekind sums,''
\emph{International Mathematics Research Notices} (2024).
arXiv:2301.00441.

\end{thebibliography}

\end{document}

% === PHASE PHI RESULT (add to Sign Pattern section) ===
% The phase constant phi=5.28 in sgn(DeltaW(p)) ~ -sgn(cos(gamma_1*log(p)+phi))
% is now derived from first principles:
% phi = -arg(rho_1 * zeta'(rho_1)) = -1.69 (mod 2pi = 5.28)
% Computed using mpmath: zeta'(rho_1) = 0.783 + 0.125i
% This confirms the explicit formula predicts BOTH frequency AND phase
% of the Chebyshev bias oscillation.
